% ==========================================================================
%  Chapter 54 — Observer/Recorder Architecture and ModelBuilder Design
% ==========================================================================

\chapter{Observer/Recorder Architecture and ModelBuilder Design}
\label{ch:observer-model-builder-design}

This chapter documents the architectural design of two major
infrastructure subsystems for the \texttt{fall\_n} library:
\begin{enumerate}
  \item an \textbf{Observer/Recorder} framework for structured monitoring,
        data extraction, and time-history recording during analyses;
  \item a \textbf{ModelBuilder} orchestrator for declarative model
        construction, absorbing the ad-hoc helpers currently scattered
        in example code.
\end{enumerate}
Both subsystems target a first stable release and adopt modern C++23
idioms (\texttt{std::print}, \texttt{std::format}) throughout.


% ══════════════════════════════════════════════════════════════════════════
\section{Motivation}
\label{sec:obs-motivation}

Three successive debugging sessions on the five-story corotational
building example (Chapter~\ref{ch:corotational-dynamic-building})
exposed a common pattern: \emph{data that should have been immediately
visible required manual inspection of internal data structures}.

\begin{itemize}
  \item The consistent mass matrix was assembled without the
        cross-sectional area (beams) or thickness (shells), producing
        dimensionally incorrect mass.  Only excessive displacements
        flagged the problem---far downstream.
  \item Gravity forces were never computed for beam and column elements
        after the shell loads were commented out.  The structure appeared
        immobile because $\mathbf{f}_{\text{gravity}} = \mathbf{0}$.
  \item All diagnostic output was embedded in a single monolithic lambda
        captured as \texttt{MonitorCallback}.  Adding a new recorder
        (e.g.\ fiber strain time-histories) required editing the lambda
        and recompiling.
\end{itemize}

These issues motivate two orthogonal improvements:
\begin{description}
  \item[Output/Monitoring] Replace the monolithic monitor callback with a
        composable \emph{observer} pipeline whose individual components
        can be attached, removed, or reconfigured at runtime without
        touching the analysis driver.
  \item[Input/Construction] Extract the $\sim$20 helper functions
        currently living in the anonymous namespace of \texttt{main.cpp}
        into a reusable, declarative \texttt{ModelBuilder} that
        orchestrates the domain--material--element--solver construction
        sequence.
\end{description}


% ══════════════════════════════════════════════════════════════════════════
\section{Current Infrastructure Audit}
\label{sec:obs-audit}

% ------------------------------------------------------------------
\subsection{Analysis Monitor Interface}

\texttt{DynamicAnalysis} defines:
\begin{lstlisting}[language=C++,basicstyle=\ttfamily\small]
using MonitorCallback =
    std::function<void(PetscInt step, double t, Vec u, Vec v)>;
void set_monitor(MonitorCallback mc);
\end{lstlisting}
A single callback is invoked at every time step.  The callback currently:
\begin{enumerate}\itemsep0pt
  \item prints progress to \texttt{stdout} (every 10 steps),
  \item computes $\max|u_x|$ and $\max|u_z|$ over all nodes,
  \item writes VTK snapshots via \texttt{StructuralVTMExporter}
        (every 100 steps),
  \item accumulates entries in a \texttt{PVDWriter}.
\end{enumerate}
All four concerns live in one lambda.  Adding a fifth (e.g.\ recording
a node time-history) requires modifying the lambda and recompiling.

% ------------------------------------------------------------------
\subsection{VTK Export}

Two output facilities exist:
\begin{itemize}
  \item \textbf{PVDWriter}: writes a ParaView \texttt{.pvd} collection
        file for time-series animation.
  \item \textbf{StructuralVTMExporter}: produces a multi-block
        \texttt{.vtm} file with five VTK sub-blocks (axis, surface,
        sections, material sites, fibers).
\end{itemize}
Neither is integrated into an observer framework; both are called
explicitly from the monitor lambda.

% ------------------------------------------------------------------
\subsection{Section Constitutive Snapshot}

The \texttt{SectionConstitutiveSnapshot} (already in the library) provides
the data bridge for recursive observation:
\begin{lstlisting}[language=C++,basicstyle=\ttfamily\small]
struct SectionConstitutiveSnapshot {
    std::optional<BeamSectionConstitutiveSnapshot>  beam;
    std::optional<ShellSectionConstitutiveSnapshot> shell;
    std::span<const FiberSectionSample>             fibers;
};
\end{lstlisting}
This gives access to per-section data (moduli, area, thickness) and
per-fiber data (position, area, strain, stress) without coupling the
observer to concrete material types.

% ------------------------------------------------------------------
\subsection{ModelBuilder (Dead Code)}

\texttt{src/model/ModelBuilder.hh} exists but is entirely commented out.
Only a default constructor and destructor compile.  The original intent
was linking \texttt{Domain} and \texttt{Model} DOF layouts, but the
design was abandoned before the current material type-erasure
architecture was in place.

% ------------------------------------------------------------------
\subsection{Helpers in \texttt{main.cpp}}

The anonymous namespace of \texttt{main.cpp} contains $\sim$20 functions
that are \emph{de facto} library API:

\begin{center}
\small
\begin{tabular}{lll}
\toprule
\textbf{Category} & \textbf{Functions} & \textbf{Target location} \\
\midrule
Section geometry &
  \texttt{rectangular\_torsion\_constant},
  \texttt{bar\_area} &
  \texttt{utils/SectionGeometry.hh} \\
Material property &
  \texttt{concrete\_initial\_modulus},
  \texttt{isotropic\_shear\_modulus} &
  \texttt{materials/MaterialProperties.hh} \\
Fiber section builder &
  \texttt{make\_steel\_fiber\_material}, &
  \texttt{materials/FiberSectionFactory.hh} \\
  & \texttt{make\_*\_concrete\_fiber\_material}, & \\
  & \texttt{add\_patch\_fibers},
    \texttt{add\_rebar\_fibers} & \\
RC section builder &
  \texttt{make\_column\_material},
  \texttt{make\_beam\_material} &
  \texttt{materials/RCSectionBuilder.hh} \\
Slab material &
  \texttt{make\_slab\_material} &
  \texttt{materials/ShellMaterialFactory.hh} \\
Domain builder &
  \texttt{add\_regular\_building\_nodes}, &
  \texttt{model/BuildingDomainBuilder.hh} \\
  & \texttt{add\_columns},
    \texttt{add\_beams},
    \texttt{add\_slabs} & \\
State queries &
  \texttt{nodal\_translation},
  \texttt{max\_translation\_norm}, &
  \texttt{post-processing/StateQuery.hh} \\
  & \texttt{max\_component\_abs} & \\
Element wrapping &
  \texttt{build\_structural\_elements} &
  \texttt{model/StructuralModelBuilder.hh} \\
VTK export &
  \texttt{export\_structural\_vtm} & (already in VTK module) \\
\bottomrule
\end{tabular}
\end{center}


% ══════════════════════════════════════════════════════════════════════════
\section{Observer/Recorder Architecture}
\label{sec:obs-architecture}

% ------------------------------------------------------------------
\subsection{Design Principles}

\begin{enumerate}
  \item \textbf{Open/Closed}: Adding a new observer must not require
        modifying \texttt{DynamicAnalysis} or any existing observer.
  \item \textbf{Single Responsibility}: Each observer handles one
        concern (progress, VTK, node recording, section recording,
        \ldots).
  \item \textbf{Composability}: A composite observer delegates to a
        vector of children, enabling rich pipelines from simple parts.
  \item \textbf{Non-intrusive}: The analysis exposes data through
        existing accessors (\texttt{Model::state\_vector()},
        \texttt{SectionConstitutiveSnapshot}).  No new ``publish'' calls
        in element code.
  \item \textbf{Efficiency}: Observers run once per time step (not per
        quadrature point).  Virtual dispatch overhead is negligible
        compared to the constitutive integration cost.
\end{enumerate}

% ------------------------------------------------------------------
\subsection{Core API}

\begin{lstlisting}[language=C++,basicstyle=\ttfamily\small]
namespace fall_n {

// Base class: one virtual call per step, plus lifecycle hooks.
template <typename ModelT>
class AnalysisObserver {
public:
    virtual ~AnalysisObserver() = default;

    // Called once before the first step.
    virtual void on_analysis_start(const ModelT& model) {}

    // Called at every time step.
    virtual void on_step(PetscInt step, double t,
                         Vec u, Vec v,
                         const ModelT& model) = 0;

    // Called once after the last step.
    virtual void on_analysis_end(const ModelT& model) {}
};

} // namespace fall_n
\end{lstlisting}

The key design choice is passing the \texttt{model} reference into
\texttt{on\_step}.  The current \texttt{MonitorCallback} only provides
\texttt{(step, t, u, v)}, forcing the lambda to capture the model.
Adding the model reference makes observers self-contained and enables
them to traverse the element--section--fiber hierarchy via
\texttt{SectionConstitutiveSnapshot} without external coupling.

% ------------------------------------------------------------------
\subsection{Composite Observer}

\begin{lstlisting}[language=C++,basicstyle=\ttfamily\small]
template <typename ModelT>
class CompositeObserver : public AnalysisObserver<ModelT> {
    std::vector<std::unique_ptr<AnalysisObserver<ModelT>>> children_;
public:
    template <typename ObsT, typename... Args>
    ObsT& add(Args&&... args) {
        auto p = std::make_unique<ObsT>(std::forward<Args>(args)...);
        auto& ref = *p;
        children_.push_back(std::move(p));
        return ref;
    }

    void on_analysis_start(const ModelT& m) override {
        for (auto& c : children_) c->on_analysis_start(m);
    }

    void on_step(PetscInt step, double t,
                 Vec u, Vec v,
                 const ModelT& m) override {
        for (auto& c : children_) c->on_step(step, t, u, v, m);
    }

    void on_analysis_end(const ModelT& m) override {
        for (auto& c : children_) c->on_analysis_end(m);
    }
};
\end{lstlisting}

% ------------------------------------------------------------------
\subsection{Integration with \texttt{DynamicAnalysis}}

The analysis class gains a new overload:
\begin{lstlisting}[language=C++,basicstyle=\ttfamily\small]
void set_observer(std::shared_ptr<AnalysisObserver<ModelT>> obs);
\end{lstlisting}
Internally, the PETSc \texttt{TSMonitor} callback delegates to
\texttt{obs->on\_step(...)}.  The old \texttt{set\_monitor(callback)}
remains for backward compatibility, wrapping the callback in a
\texttt{LambdaObserver} adapter.

% ------------------------------------------------------------------
\subsection{Concrete Observers}

\subsubsection{ConsoleProgressObserver}

Replaces the ``print every $N$ steps'' logic.  Parameters: print
interval, component selection.

\begin{lstlisting}[language=C++,basicstyle=\ttfamily\small]
template <typename ModelT>
class ConsoleProgressObserver : public AnalysisObserver<ModelT> {
    int interval_;
public:
    explicit ConsoleProgressObserver(int interval = 10)
        : interval_{interval} {}

    void on_step(PetscInt step, double t,
                 Vec u, Vec v,
                 const ModelT& model) override {
        if (step % interval_ != 0) return;
        const double max_u = max_translation_norm(model);
        std::println("  [step {:4d}]  t = {:8.4f} s   max|u| = {:10.6e} m",
                     static_cast<int>(step), t, max_u);
    }
};
\end{lstlisting}

\subsubsection{VTKSnapshotObserver}

Encapsulates the \texttt{StructuralVTMExporter} + \texttt{PVDWriter}
pipeline.  Parameters: output directory, snapshot interval, exporter
configuration.

\begin{lstlisting}[language=C++,basicstyle=\ttfamily\small]
template <typename ModelT, typename BeamProfileT, typename ThicknessProfileT>
class VTKSnapshotObserver : public AnalysisObserver<ModelT> {
    std::string output_dir_;
    int interval_;
    PVDWriter pvd_;
    BeamProfileT beam_profile_;
    ThicknessProfileT thickness_profile_;
public:
    // constructor omitted for brevity

    void on_step(PetscInt step, double t, Vec u, Vec v,
                 const ModelT& model) override {
        if (step % interval_ != 0) return;
        auto filename = std::format("{}/snapshot_{:06d}.vtm",
                                    output_dir_, static_cast<int>(step));
        fall_n::vtk::StructuralVTMExporter exporter(
            model, beam_profile_, thickness_profile_);
        exporter.write(filename);
        pvd_.add_timestep(t, filename);
    }

    void on_analysis_end(const ModelT&) override {
        pvd_.write();
    }
};
\end{lstlisting}

\subsubsection{NodeRecorder}

Records time-history of selected DOFs at specific nodes.  Stores data
in columnar format for post-processing or CSV export.

\begin{lstlisting}[language=C++,basicstyle=\ttfamily\small]
template <typename ModelT>
class NodeRecorder : public AnalysisObserver<ModelT> {
    struct Channel {
        std::size_t node_id;
        std::size_t dof;          // local DOF index (0..5)
        std::vector<double> data; // time-history samples
    };
    std::vector<Channel> channels_;
    std::vector<double>  times_;
    int interval_;
public:
    void record(std::size_t node, std::size_t dof);

    void on_step(PetscInt step, double t, Vec u, Vec v,
                 const ModelT& model) override;

    // Accessor: returns span of recorded values for a channel
    std::span<const double> history(std::size_t channel) const;
    std::span<const double> time_axis() const;

    // Export to CSV
    void write_csv(const std::string& filename) const;
};
\end{lstlisting}

\subsubsection{SectionRecorder}

Records section-level quantities (e.g.\ peak fiber strain, moment
resultant) by traversing the element--section hierarchy via
\texttt{SectionConstitutiveSnapshot}.

\subsubsection{EnergyRecorder}

Records kinetic, potential (strain), and dissipated energy per step.
Useful for verifying energy balance in nonlinear dynamics.


% ------------------------------------------------------------------
\subsection{Data Flow Diagram}

\begin{center}
\begin{minipage}{0.80\textwidth}
\small
\begin{verbatim}
DynamicAnalysis
  |-- solve(): on_analysis_start(model)
  |-- TSMonitor callback (PETSc)
       |
       v
  ObserverCallback<ModelT>    (type-erased bridge)
       |
       +-- [static path]  CompositeObserver<ModelT, Obs...>
       |       tuple<Obs...> + fold expressions (zero vtable)
       |
       +-- [dynamic path]  DynamicObserverList<ModelT>
               vector<unique_ptr<ObserverBase<ModelT>>>
       |
       |-- ConsoleProgressObserver  --> stdout (std::println)
       |-- VTKSnapshotObserver      --> .vtm/.pvd files
       |-- NodeRecorder             --> in-memory + CSV
       |-- ElementRecorder          --> in-memory (TODO: accessor)
       |-- FiberRecorder            --> in-memory (TODO: accessor)
       |-- EnergyRecorder           --> in-memory (v^T M v)
       |-- MaxResponseTracker       --> peak envelope
       |
       +-- (user-defined observers)
  |
  |-- solve(): on_analysis_end(model)
\end{verbatim}
\end{minipage}
\end{center}

Each observer independently reads from the \texttt{Model} via its
existing accessors.  No observer modifies the model state.


% ------------------------------------------------------------------
\subsection{Design Alternatives Considered}

\paragraph{Type-erased observer (SBO).}
The \texttt{Material<T>} pattern uses small-buffer optimisation for
value semantics.  For observers this is unnecessary: observers are
few (typically 3--5), created once, and stored by pointer.  Virtual
dispatch through \texttt{unique\_ptr} is simpler and sufficient.

\paragraph{Concept-based observer.}
A C++20 concept (\texttt{AnalysisObserverConcept}) could eliminate the
vtable.  However, the composite pattern requires homogeneous storage
(\texttt{vector<unique\_ptr<Base>>}), which reintroduces virtual
dispatch.  The concept adds complexity without benefit here.

\paragraph{Signal/slot (e.g.\ \texttt{boost::signals2}).}
Decouples publishers and subscribers at the cost of an external
dependency.  The library avoids non-essential dependencies; the
composite pattern achieves the same decoupling.

\paragraph{Callback vector.}
\texttt{DynamicAnalysis} could store a
\texttt{vector<MonitorCallback>} instead of a single callback.
This is simpler but loses lifecycle hooks (\texttt{on\_analysis\_start},
\texttt{on\_analysis\_end}) and named semantics.  The observer base
class is a small price for the richer protocol.


% ══════════════════════════════════════════════════════════════════════════
\section{Comparison with OpenSees Recorder Architecture}
\label{sec:opensees-comparison}

The OpenSees framework~\cite{OpenSees}%
\footnote{F.~McKenna, G.\,L.~Fenves, M.\,H.~Scott, \emph{Open System
for Earthquake Engineering Simulation}, 2000.  Source:
\url{https://github.com/OpenSees/OpenSees}.}
provides a \texttt{Recorder} subsystem that superficially resembles
the Observer/Recorder pipeline described in this chapter.
This section analyses the two architectures---OpenSees' runtime
scripting--driven \texttt{Recorder} and \texttt{fall\_n}'s
compile-time polymorphic Observer---highlighting points of convergence
and deliberate divergence.

% ------------------------------------------------------------------
\subsection{OpenSees Recorder: Key Architectural Traits}
\label{sec:opensees-recorder}

In OpenSees, a \texttt{Recorder} is a polymorphic object derived from
an abstract base class with the following interface skeleton:

\begin{lstlisting}[language=C++,basicstyle=\small\ttfamily]
class Recorder : public MovableObject, public TaggedObject {
  virtual int record(int commitTag, double timeStamp) = 0;
  virtual int restart()                               = 0;
  virtual int setDomain(Domain& theDomain)            = 0;
  // ... serialization, copy, class-tag machinery
};
\end{lstlisting}

Concrete implementations include \texttt{NodeRecorder},
\texttt{ElementRecorder}, \texttt{EnvelopeNodeRecorder},
\texttt{DriftRecorder}, \texttt{PatternRecorder}, among others.
Notable architectural characteristics:

\begin{enumerate}
  \item \textbf{Domain coupling.}  Every recorder is explicitly
        \emph{set} on a \texttt{Domain} via \texttt{setDomain()}.
        The recorder then queries the domain's node/element maps
        internally.  This couples the recorder tightly to the
        \texttt{Domain} class hierarchy.

  \item \textbf{File-oriented output.}  Most concrete recorders
        write to disk (\texttt{FileStream}, \texttt{DataFileStream},
        \texttt{XmlFileStream}).  The ``record'' verb implies
        \emph{persistence}---data goes to a file, not into an
        in-memory buffer that later observers can consume.

  \item \textbf{Tag-based identification.}  Recorders inherit from
        \texttt{TaggedObject}, receiving an integer tag that the Tcl
        (or Python) wrapper resolves.  The scripting layer is the
        primary configuration mechanism:
        \texttt{recorder Node -file out.txt -node 1 2 3 -dof 1 disp}.

  \item \textbf{Serialisation infrastructure.}  Through
        \texttt{MovableObject}, recorders participate in distributed
        analysis and checkpoint/restart.  This adds a non-trivial
        interface burden (\texttt{sendSelf}, \texttt{recvSelf},
        \texttt{classTag}).

  \item \textbf{Pull model.}  The analysis driver calls
        \texttt{recorder->record(commitTag, timeStamp)} at each
        converged step; the recorder \emph{pulls} data from the
        domain by navigating its node/element maps.
\end{enumerate}

% ------------------------------------------------------------------
\subsection{\texttt{fall\_n} Observer: Contrasting Design}
\label{sec:falln-observer-contrast}

The \texttt{fall\_n} observer pipeline differs in several
architectural dimensions:

\begin{enumerate}
  \item \textbf{Push model.}  The analysis hands a
        \texttt{StepEvent} value (containing time, step number, and
        a \emph{const reference} to the model) to each observer.
        Observers do not query a global domain; they receive exactly
        the data they need.  This eliminates the
        \texttt{setDomain()} coupling entirely.

  \item \textbf{In-memory first.}  Observers like
        \texttt{NodeRecorder} accumulate data in
        \texttt{std::vector} time-series.  File export is a
        \emph{secondary step} (e.g.\ \texttt{write\_csv()}).
        This permits post-analysis querying without re-reading files.

  \item \textbf{Lifecycle protocol.}
        \texttt{on\_analysis\_start}, \texttt{on\_step},
        \texttt{on\_analysis\_end}---a richer contract than OpenSees'
        single \texttt{record()} entry point.  Start/end hooks enable
        resource acquisition (opening files, allocating buffers) and
        clean teardown (flushing, summary statistics) without
        sentinel logic inside \texttt{on\_step}.

  \item \textbf{No scripting layer dependency.}  Configuration is
        entirely C++ (compile-time checked).  Observers are added via
        the \texttt{CompositeObserver} builder or directly on the
        analysis object.  No tag system, no class-tag registry, no
        Tcl interpreter.

  \item \textbf{Composability.}  \texttt{CompositeObserver} stores a
        \texttt{vector<unique\_ptr<AnalysisObserver>>}, dispatching
        each hook to all children.  OpenSees achieves the same
        multiplicity through the \texttt{Domain}'s recorder list,
        but the composition is implicit (the domain iterates its
        recorders) rather than explicit (a first-class composite
        object).

  \item \textbf{No serialisation burden.}  Observers have no
        obligation to implement \texttt{sendSelf/recvSelf}.  In a
        future MPI-distributed analysis, serialisation can be added
        only to the observers that need it, not imposed on all via
        a base class.

  \item \textbf{Type safety.}  The observer receives a
        \texttt{const Model\&}, from which it can call
        \texttt{node\_displacement()}, \texttt{node\_velocity()}, etc.
        OpenSees recorders downcast through \texttt{Domain} to
        obtain \texttt{Node*} pointers, then call
        \texttt{getDisp()} --- a less type-safe chain.
\end{enumerate}

% ------------------------------------------------------------------
\subsection{Architectural Summary Table}
\label{sec:comparison-table}

\begin{center}
\small
\begin{tabular}{p{3.8cm}p{5cm}p{5cm}}
\toprule
\textbf{Aspect} & \textbf{OpenSees Recorder} & \textbf{\texttt{fall\_n} Observer} \\
\midrule
Data flow
  & Pull: recorder queries Domain
  & Push: analysis delivers StepEvent \\
\addlinespace
Primary output
  & File streams (disk-first)
  & In-memory vectors (export later) \\
\addlinespace
Lifecycle hooks
  & \texttt{record()} only
  & \texttt{on\_start}, \texttt{on\_step}, \texttt{on\_end} \\
\addlinespace
Composition
  & Implicit (Domain iterates list)
  & Explicit (\texttt{CompositeObserver}) \\
\addlinespace
Configuration
  & Tcl/Python scripting + class tags
  & Compile-time C++ \\
\addlinespace
Domain coupling
  & \texttt{setDomain()} required
  & None (const model ref in event) \\
\addlinespace
Serialisation
  & Mandatory (\texttt{MovableObject})
  & Optional (none by default) \\
\addlinespace
Base class weight
  & Heavy (\texttt{MovableObject} + \texttt{TaggedObject})
  & Minimal (single pure-virtual base) \\
\addlinespace
Type safety
  & \texttt{void*} / downcast chains
  & Strongly typed model queries \\
\bottomrule
\end{tabular}
\end{center}

% ------------------------------------------------------------------
\subsection{Shared Conceptual Ground}
\label{sec:comparison-shared}

Despite the architectural divergence, several concepts are shared:

\begin{itemize}
  \item \textbf{Observer pattern at core.}  Both systems implement
        the GoF Observer pattern: the analysis is the subject;
        recorders/observers are subscribers notified at each
        converged step.

  \item \textbf{Concrete specialisation by entity type.}
        \texttt{NodeRecorder} and \texttt{ElementRecorder} exist in
        both systems, reflecting the natural partition of response
        quantities.

  \item \textbf{Open extension.}  Users can (and do) add new recorder
        types in OpenSees by subclassing \texttt{Recorder}.  In
        \texttt{fall\_n}, the same is achieved by subclassing
        \texttt{AnalysisObserver}.

  \item \textbf{Per-step granularity.}  Both systems operate at
        the converged-step level: one notification per accepted
        time/load step, with access to the current state.

  \item \textbf{Separation of analysis logic from output.}  Neither
        system embeds output formatting inside the solver.
        The analysis code is agnostic to what observers/recorders do
        with the data.
\end{itemize}

\paragraph{Design philosophy.}
The key difference is one of \emph{architectural epoch}.  OpenSees was
designed in the late 1990s for a Tcl scripting environment, emphasising
runtime flexibility and file persistence.  \texttt{fall\_n} is designed
for modern C++ (C++20/23), emphasising compile-time safety, value
semantics, and in-memory composability.  The Observer \emph{pattern}
is the same; the \emph{realisation} reflects two decades of language
and design evolution.


% ══════════════════════════════════════════════════════════════════════════
\section{ModelBuilder Architecture}
\label{sec:model-builder}

% ------------------------------------------------------------------
\subsection{Layered Design}

The ModelBuilder is split into two layers:

\begin{description}
  \item[Layer 1: \texttt{ModelBuilder<dim>}] (generic) Handles domain
        construction, DOF configuration, constraint application, and
        analysis setup---independent of element type.
  \item[Layer 2: \texttt{StructuralModelBuilder}] (specialised) Extends
        Layer~1 with frame/shell convenience methods: material-to-group
        mapping, fiber section builders, VTK export configuration, and
        structural element wrapping.
\end{description}

This separation allows future continuum-only problems to use Layer~1
alone, while structural engineering workflows get the full convenience
of Layer~2.

% ------------------------------------------------------------------
\subsection{Layer 1: Generic ModelBuilder}

\begin{lstlisting}[language=C++,basicstyle=\ttfamily\small]
template <std::size_t dim>
class ModelBuilder {
public:
    // Domain construction
    ModelBuilder& add_node(PetscInt id, double x, double y, double z = 0.0);
    ModelBuilder& add_grid(std::span<const double> x,
                           std::span<const double> y,
                           std::span<const double> z);

    template <typename ElemType, typename IntStrategy>
    ModelBuilder& add_element(std::size_t tag,
                              std::span<const PetscInt> connectivity,
                              std::string_view group = "");

    ModelBuilder& assemble();

    // Constraints
    ModelBuilder& fix_base(double z = 0.0);
    ModelBuilder& fix_node(PetscInt id);
    ModelBuilder& constrain_dof(PetscInt node, std::size_t dof, double val);

    // Boundary conditions
    ModelBuilder& add_force(PetscInt node,
                            std::array<double, dim> components,
                            time_fn::TimeFunction tf = time_fn::constant(1.0));
    ModelBuilder& add_ground_motion(std::size_t dir, /* ... */);

    // Finalize
    Domain<dim>& domain();

private:
    Domain<dim> domain_;
};
\end{lstlisting}

The fluent (method-chaining) API allows declarative construction:
\begin{lstlisting}[language=C++,basicstyle=\ttfamily\small]
auto builder = ModelBuilder<3>{}
    .add_grid(X_GRID, Y_GRID, Z_LEVELS)
    .fix_base(0.0)
    .assemble();
\end{lstlisting}

% ------------------------------------------------------------------
\subsection{Layer 2: StructuralModelBuilder}

\begin{lstlisting}[language=C++,basicstyle=\ttfamily\small]
template <typename FrameElemT, typename ShellElemT>
class StructuralModelBuilder {
    using StructuralPolicy = SingleElementPolicy<StructuralElement>;
    using ModelT = Model</* ... */, StructuralPolicy>;
public:
    // Material assignment by physical group
    StructuralModelBuilder& set_frame_material(
        std::string_view group,
        const Material<TimoshenkoBeam3D>& mat);

    StructuralModelBuilder& set_shell_material(
        std::string_view group,
        const Material<MindlinReissnerShell3D>& mat);

    // Convenience: RC fiber section builders
    StructuralModelBuilder& set_rc_column_section(
        std::string_view group,
        const RCColumnSectionSpec& spec);

    StructuralModelBuilder& set_rc_beam_section(
        std::string_view group,
        const RCBeamSectionSpec& spec);

    // Gravity
    StructuralModelBuilder& apply_self_weight(double density, double g);

    // Build
    ModelT build(Domain<3>& domain);

private:
    std::map<std::string, Material<TimoshenkoBeam3D>> frame_materials_;
    std::map<std::string, Material<MindlinReissnerShell3D>> shell_materials_;
};
\end{lstlisting}

The \texttt{RCColumnSectionSpec} / \texttt{RCBeamSectionSpec} structs
absorb the parameters currently hardcoded in the \texttt{make\_column\_material()}
and \texttt{make\_beam\_material()} helpers.

% ------------------------------------------------------------------
\subsection{Framed Structure Convenience}

For regular framed buildings (the most common use case), a higher-level
builder provides a grid-based API:

\begin{lstlisting}[language=C++,basicstyle=\ttfamily\small]
namespace fall_n::building {

struct FramedBuildingSpec {
    std::vector<double> x_axes;
    std::vector<double> y_axes;
    int num_stories;
    double story_height;
    // per-group section specs, slab thickness, etc.
};

auto make_framed_building(const FramedBuildingSpec& spec)
    -> std::pair<Domain<3>, StructuralModelBuilder::ModelT>;

} // namespace fall_n::building
\end{lstlisting}

This replaces the \texttt{add\_regular\_building\_nodes()},
\texttt{add\_columns()}, \texttt{add\_beams()}, \texttt{add\_slabs()}
functions, reducing the entire construction sequence to a single call.


% ------------------------------------------------------------------
\subsection{Helper Function Extraction Plan}

Functions currently in \texttt{main.cpp} are migrated as follows:

\paragraph{Phase A --- Pure utilities (no dependencies).}
\begin{itemize}
  \item \texttt{rectangular\_torsion\_constant(b, h)}
        $\to$ \texttt{src/utils/SectionProperties.hh}
  \item \texttt{concrete\_initial\_modulus(fpc)}
        $\to$ \texttt{src/materials/ConcreteMaterial.hh}
  \item \texttt{isotropic\_shear\_modulus(E, nu)}
        $\to$ \texttt{src/materials/MaterialProperties.hh}
  \item \texttt{bar\_area(diameter)}
        $\to$ \texttt{src/utils/SectionProperties.hh}
\end{itemize}

\paragraph{Phase B --- Fiber section factories.}
\begin{itemize}
  \item \texttt{add\_patch\_fibers()}
        $\to$ \texttt{src/materials/FiberSectionFactory.hh}
  \item \texttt{add\_rebar\_fibers()}
        $\to$ same header
  \item \texttt{make\_steel\_fiber\_material()},
        \texttt{make\_*\_concrete\_fiber\_material()}
        $\to$ \texttt{src/materials/FiberMaterialFactory.hh}
  \item \texttt{make\_column\_material()},
        \texttt{make\_beam\_material()}
        $\to$ \texttt{src/materials/RCSectionBuilder.hh}
\end{itemize}

\paragraph{Phase C --- State queries.}
\begin{itemize}
  \item \texttt{nodal\_translation()},
        \texttt{max\_translation\_norm()},
        \texttt{max\_component\_abs()}
        $\to$ \texttt{src/post-processing/StateQuery.hh}
\end{itemize}

\paragraph{Phase D --- Domain+model builders.}
\begin{itemize}
  \item \texttt{add\_regular\_building\_nodes()}, \texttt{add\_columns()},
        \texttt{add\_beams()}, \texttt{add\_slabs()}
        $\to$ \texttt{src/model/BuildingDomainBuilder.hh}
  \item \texttt{build\_structural\_elements()}
        $\to$ \texttt{src/model/StructuralModelBuilder.hh}
\end{itemize}


% ══════════════════════════════════════════════════════════════════════════
\section{Modern C++ Migration}
\label{sec:cpp23-migration}

% ------------------------------------------------------------------
\subsection{std::print and std::format}

The library already enables C++23 (\texttt{-std=c++2b}) and includes
\texttt{<format>} and \texttt{<print>} in several headers, but most
output still uses \texttt{std::cout}.  The migration:

\begin{itemize}
  \item Replace \texttt{std::cout << "text" << value << "\textbackslash n"}
        with \texttt{std::println("text \{\}", value)}.
  \item Replace \texttt{std::ostringstream} formatting with
        \texttt{std::format}.
  \item Prioritise new code (observers, builders) for
        \texttt{std::print}; migrate existing code incrementally.
\end{itemize}

\paragraph{Example.}
Current:
\begin{lstlisting}[language=C++,basicstyle=\ttfamily\small]
std::cout << "  Density (RC)       : " << RC_DENSITY << " MN s^2/m^4\n";
\end{lstlisting}
Becomes:
\begin{lstlisting}[language=C++,basicstyle=\ttfamily\small]
std::println("  Density (RC)       : {} MN s^2/m^4", RC_DENSITY);
\end{lstlisting}


% ══════════════════════════════════════════════════════════════════════════
\section{Implementation Order and Prerequisites}
\label{sec:obs-impl-order}

% ------------------------------------------------------------------
\subsection{Prerequisites}

Before beginning implementation:
\begin{enumerate}
  \item \textbf{Test suite health}: The two pre-existing failures
        (\texttt{cantilever\_benchmark} and
        \texttt{snes\_gmsh\_pipeline}, both dependent on external mesh
        files) were removed from \texttt{CMakeLists.txt}.
        Suite now stands at 45/45. \emph{Verified.}
  \item \textbf{Git commit}: The current working state (mass-matrix
        fix + force-application fix) must be committed to preserve a
        known-good baseline.
  \item \textbf{Build clean}: The build must compile with zero
        warnings under \texttt{-Werror}.
        \emph{Verified.}
\end{enumerate}

% ------------------------------------------------------------------
\subsection{Phased Implementation}

\begin{description}
  \item[Phase 1: Observer framework] (this chapter)
    \begin{enumerate}
      \item Define \texttt{AnalysisObserver<ModelT>} base class.
      \item Implement \texttt{CompositeObserver}.
      \item Implement \texttt{ConsoleProgressObserver}.
      \item Implement \texttt{VTKSnapshotObserver}.
      \item Integrate into \texttt{DynamicAnalysis}
            (add \texttt{set\_observer()}, keep backward-compatible
            \texttt{set\_monitor()}).
      \item Refactor \texttt{main.cpp} to use the composite pipeline.
      \item Add \texttt{NodeRecorder} and verify with the building example.
    \end{enumerate}

  \item[Phase 2: Helper extraction]
    \begin{enumerate}
      \item Extract pure utilities (Phase~A from
            \S\ref{sec:model-builder}).
      \item Extract fiber section factories (Phase~B).
      \item Extract state queries (Phase~C).
      \item Verify \texttt{main.cpp} compiles using the new headers.
    \end{enumerate}

  \item[Phase 3: ModelBuilder]
    \begin{enumerate}
      \item Implement \texttt{ModelBuilder<dim>} (Layer~1).
      \item Implement \texttt{StructuralModelBuilder} (Layer~2).
      \item Implement \texttt{make\_framed\_building()} convenience.
      \item Refactor \texttt{main.cpp} to use the builders.
      \item Delete dead code in \texttt{ModelBuilder.hh}.
    \end{enumerate}

  \item[Phase 4: std::print migration]
    \begin{enumerate}
      \item Migrate all new code (observers, builders) to
            \texttt{std::print}.
      \item Incrementally migrate existing library output.
    \end{enumerate}
\end{description}


% ══════════════════════════════════════════════════════════════════════════
\section{Trade-offs and Design Decisions}
\label{sec:obs-tradeoffs}

\begin{center}
\small
\begin{tabular}{p{4cm}p{4.5cm}p{4.5cm}}
\toprule
\textbf{Decision} & \textbf{Pros} & \textbf{Cons} \\
\midrule
Virtual base class for observers &
  Simple, familiar, composable via \texttt{unique\_ptr} &
  Small vtable overhead (negligible for per-step calls) \\
\addlinespace
Template \texttt{ModelT} on observer &
  Works with any model type without casting &
  Observers are not exchangeable across model types \\
\addlinespace
Fluent (chaining) API for builder &
  Readable, declarative construction &
  Cannot easily validate intermediate states \\
\addlinespace
Two-layer builder &
  Generic layer reusable for continuum; structural layer adds convenience &
  More classes to maintain \\
\addlinespace
Pass \texttt{model} to \texttt{on\_step} &
  Observer is self-contained; no captured references &
  Slight API change vs current \texttt{MonitorCallback} \\
\addlinespace
\texttt{shared\_ptr} for observer in analysis &
  Observer lifetime may extend beyond analysis (e.g.\ post-processing) &
  Shared ownership adds minor complexity \\
\bottomrule
\end{tabular}
\end{center}


% ══════════════════════════════════════════════════════════════════════════
\section{Phase~1 Implementation Notes}
\label{sec:obs-implementation}

This section records the actual implementation choices, deviations from
the design proposal, and the file organisation that resulted from
Phase~1 development.

% ------------------------------------------------------------------
\subsection{Dual-flavour Observer Composition}
\label{sec:impl-dual-flavour}

The design document (\S\ref{sec:obs-architecture}) proposed a single
virtual-dispatch model using
\texttt{vector<unique\_ptr<AnalysisObserver<ModelT>{}>{}>{}>}.
During implementation a dual-flavour architecture was adopted to
honour the user requirement of \emph{favouring compile-time efficiency}:

\begin{description}
  \item[Static path --- \texttt{CompositeObserver<ModelT, Obs...>}]
    Stores observers in a \texttt{std::tuple<Obs...>}.
    Dispatch uses C++17 fold expressions:
\begin{lstlisting}[language=C++,basicstyle=\ttfamily\small]
void on_step(const StepEvent& ev, const ModelT& model) {
    std::apply([&](auto&... obs) {
        (obs.on_step(ev, model), ...);
    }, observers_);
}
\end{lstlisting}
    This eliminates \emph{all} virtual dispatch overhead: no vtable
    lookup, no indirect branch prediction miss, and full inlining
    opportunity for trivial observers such as
    \texttt{ConsoleProgressObserver}.  The factory
    \texttt{make\_composite\_observer<ModelT>(obs...)} deduces types
    automatically.

  \item[Dynamic path --- \texttt{DynamicObserverList<ModelT>}]
    Stores observers via
    \texttt{vector<unique\_ptr<ObserverBase<ModelT>{}>{}>{}>}
    for runtime composition (e.g.\ user
    scripts, configuration-file--driven pipelines).  Preserves the
    classical composite pattern from the original design.
\end{description}

Both flavours satisfy the \texttt{ObserverBase<ModelT>} concept (virtual
base class) but the static path can bypass it entirely when composed via
\texttt{CompositeObserver}.

% ------------------------------------------------------------------
\subsection{StepEvent Value Type}
\label{sec:impl-step-event}

Instead of passing four arguments (\texttt{step, t, u, v}) to
\texttt{on\_step}, the implementation introduces a lightweight
aggregate:

\begin{lstlisting}[language=C++,basicstyle=\ttfamily\small]
struct StepEvent {
    PetscInt step;
    double   time;
    Vec      displacement;   // PETSc Vec (non-owning handle)
    Vec      velocity;       // PETSc Vec (non-owning handle)
};
\end{lstlisting}

This simplifies signatures, is trivially copyable (two scalars + two
opaque pointers), and is easily extensible (e.g.\ adding
\texttt{acceleration} for Newmark schemes) without changing every
observer signature.

% ------------------------------------------------------------------
\subsection{ObserverCallback Bridge}
\label{sec:impl-observer-callback}

\texttt{DynamicAnalysis} stores a PETSc \texttt{TSMonitor} callback
whose signature is fixed by the C library.  A type-erased bridge,
\texttt{ObserverCallback<ModelT>}, wraps any observer flavour into three
\texttt{std::function} members:

\begin{lstlisting}[language=C++,basicstyle=\ttfamily\small]
template <typename ModelT>
struct ObserverCallback {
    std::function<void(const ModelT&)>                    on_start;
    std::function<void(const StepEvent&, const ModelT&)>  on_step;
    std::function<void(const ModelT&)>                    on_end;
};
\end{lstlisting}

The factory \texttt{make\_observer\_callback(Obs\&)} constructs this
from any observer (static or dynamic) by binding references.
\texttt{DynamicAnalysis::Monitor()} feeds into \texttt{on\_step};
lifecycle hooks are called from \texttt{solve()}.  The old
\texttt{MonitorCallback} is preserved for backward compatibility:
both are invoked sequentially.

%-------------------------------------------------------------------
\subsection{Concrete Observers (Implemented)}
\label{sec:impl-concrete-observers}

Seven observers were implemented in
\texttt{src/analysis/Observers.hh}:

\begin{center}
\small
\begin{tabular}{lp{7cm}}
\toprule
\textbf{Observer} & \textbf{Description} \\
\midrule
\texttt{ConsoleProgressObserver} &
  Prints step/time/max-displacement every $N$ steps via
  \texttt{std::println}. \\
\texttt{VTKSnapshotObserver} &
  Writes VTM snapshots and maintains PVD time series via a
  user-supplied exporter factory (avoids hard VTK include). \\
\texttt{NodeRecorder} &
  Records DOF time-histories at selected (node, dof) channels;
  CSV export. \\
\texttt{ElementRecorder} &
  Per-element section resultants (axial, shear, moment).
  \emph{Requires \texttt{section\_snapshots()} accessor on
  \texttt{StructuralElement}.} \\
\texttt{FiberRecorder} &
  Per-fiber stress/strain time-histories.
  \emph{Same accessor prerequisite.} \\
\texttt{EnergyRecorder} &
  Kinetic energy $\frac{1}{2}\mathbf{v}^{\!\top}\!\mathbf{M}\mathbf{v}$
  and cumulative external work. \\
\texttt{MaxResponseTracker} &
  Tracks peak $|\mathbf{u}|$, $|u_x|$, $|u_y|$, $|u_z|$ with
  time stamps. \\
\bottomrule
\end{tabular}
\end{center}

\texttt{VTKSnapshotObserver} deserves special note: its second template
parameter \texttt{ExporterFactory} is a callable
\texttt{(const ModelT\&) -> Exporter} rather than a direct VTK type
reference.  This \emph{compile-time firewall} prevents
\texttt{Observers.hh} from pulling in VTK translation-unit headers,
keeping the observer framework header lightweight.

% ------------------------------------------------------------------
\subsection{StateQuery Extraction}
\label{sec:impl-state-query}

The helper functions \texttt{nodal\_translation()},
\texttt{max\_translation\_norm()}, and \texttt{max\_component\_abs()}
were extracted from the anonymous namespace in \texttt{main.cpp} into
\texttt{src/post-processing/StateQuery.hh} under namespace
\texttt{fall\_n::query}.  These are now reusable across any analysis
driver or observer and serve as the standard response-quantity accessors.

\texttt{main.cpp} was updated to call \texttt{fall\_n::query::} instead
of local duplicates; the local template functions were deleted.

% ------------------------------------------------------------------
\subsection{\texttt{std::print/format} Migration (Partial Phase~2)}
\label{sec:impl-print-migration}

All new observer code uses \texttt{std::println} and
\texttt{std::format} exclusively.  Additionally, the entire
\texttt{main.cpp} dynamic-analysis driver was migrated:
\begin{itemize}
  \item Dynamic parameter block $\to$ \texttt{std::println}.
  \item Final report (geometry, materials, solver outcome, response
        quantities) $\to$ \texttt{std::println} with
        \texttt{\{:.6f\}} formatting.
  \item \texttt{<iostream>} and \texttt{<iomanip>} removed from
        \texttt{main.cpp} includes.
\end{itemize}

% ------------------------------------------------------------------
\subsection{File Organisation}
\label{sec:impl-file-org}

\begin{center}
\small
\begin{tabular}{lp{8cm}}
\toprule
\textbf{File} & \textbf{Contents} \\
\midrule
\texttt{src/analysis/AnalysisObserver.hh} &
  \texttt{StepEvent}, \texttt{ObserverBase<ModelT>},
  \texttt{DynamicObserverList<ModelT>},
  \texttt{CompositeObserver<ModelT, Obs...>},
  \texttt{ObserverCallback<ModelT>},
  factory functions. \\
\texttt{src/analysis/Observers.hh} &
  All 7 concrete observer implementations +
  \texttt{make\_vtk\_observer} factory. \\
\texttt{src/post-processing/StateQuery.hh} &
  \texttt{fall\_n::query::nodal\_translation()},
  \texttt{max\_translation\_norm()},
  \texttt{max\_component\_abs()},
  \texttt{nodal\_dof\_value()}. \\
\texttt{src/analysis/DynamicAnalysis.hh} &
  Modified: \texttt{observer\_callback\_} member,
  \texttt{set\_observer()} overloads,
  lifecycle hooks in \texttt{solve()}. \\
\texttt{src/analysis/Analysis.hh} &
  Updated umbrella include. \\
\texttt{header\_files.hh} &
  Added \texttt{StateQuery.hh}. \\
\bottomrule
\end{tabular}
\end{center}

% ------------------------------------------------------------------
\subsection{DynamicAnalysis Integration Detail}
\label{sec:impl-da-integration}

The integration preserves full backward compatibility:

\begin{enumerate}
  \item \texttt{DynamicAnalysis} gains a private member
    \texttt{ObserverCallback<ModelT> observer\_callback\_}.
  \item The static \texttt{Monitor(TS, PetscInt, PetscReal, Vec, void*)}
    callback first invokes \texttt{observer\_callback\_.on\_step}
    (wrapping arguments into a \texttt{StepEvent}), then the legacy
    \texttt{monitor\_callback\_}.
  \item Two \texttt{set\_observer()} overloads:
    \begin{itemize}
      \item \texttt{set\_observer(ObserverCallback<ModelT>)} for
            pre-built callbacks;
      \item \texttt{template<class Obs> set\_observer(Obs\&)} that
            calls \texttt{make\_observer\_callback(obs)}.
    \end{itemize}
  \item Both \texttt{solve()} overloads invoke
    \texttt{on\_analysis\_start} before \texttt{TSSolve} and
    \texttt{on\_analysis\_end} after the post-processing phase.
\end{enumerate}

% ------------------------------------------------------------------
\subsection{Test Suite After Phase~1}
\label{sec:impl-test-suite}

Two pre-existing failing tests (\texttt{cantilever\_benchmark} and
\texttt{snes\_gmsh\_pipeline}) that depended on external mesh files
were removed from \texttt{CMakeLists.txt}.  After all Phase~1 changes
the test suite passes:

\begin{center}
  \textbf{45 / 45 tests passed, 0 failures.}
\end{center}

The build compiles under \texttt{clang++-20} with \texttt{-Werror}
and produces zero warnings.

% ------------------------------------------------------------------
\subsection{Known Limitations and Next Steps}
\label{sec:impl-limitations}

\begin{itemize}
  \item \texttt{ElementRecorder} and \texttt{FiberRecorder} store
        empty records because \texttt{StructuralElement}'s type-erased
        interface does not yet expose \texttt{sections()} or
        \texttt{section\_snapshots()}.  Adding this accessor to the
        \texttt{Concept} vtable is a prerequisite for Phase~2.
  \item \texttt{LinearAnalysis} and \texttt{NonlinearAnalysis} do not
        have monitor support.  Extending the observer integration to
        iterative solvers (SNES monitoring) is future work.
  \item The \texttt{ModelBuilder} (Phase~3) has not been started.
  \item Node recorder CSV export has been implemented but not yet
        exercised in a test.
\end{itemize}


% ══════════════════════════════════════════════════════════════════════════
\section{Phase~2 Implementation Notes: Helper Extraction}
\label{sec:phase2-implementation}

Phase~2 extracted the pure-utility functions, fiber-section factories,
and RC section builders from the anonymous namespace of
\texttt{main.cpp} into four new reusable library headers.
The extraction followed the plan in \S\ref{sec:model-builder}
(Phases~A through~C).

% ------------------------------------------------------------------
\subsection{New File Organisation}
\label{sec:phase2-files}

\begin{center}
\small
\begin{tabular}{lp{8cm}}
\toprule
\textbf{Header} & \textbf{Contents} \\
\midrule
\texttt{src/utils/SectionProperties.hh}
  & Pure \texttt{constexpr} utilities:
    \texttt{rectangular\_torsion\_constant()},
    \texttt{bar\_area()},
    \texttt{isotropic\_shear\_modulus()},
    \texttt{concrete\_initial\_modulus()}.
    Depends only on \texttt{<cmath>} and \texttt{<numbers>}. \\
\addlinespace
\texttt{src/materials/FiberSectionFactory.hh}
  & Generic fiber-mesh discretization helpers
    (\texttt{add\_patch\_fibers()}, \texttt{add\_rebar\_fibers()})
    and parameterised uniaxial material factories
    (\texttt{make\_steel\_fiber\_material()},
     \texttt{make\_unconfined\_concrete()},
     \texttt{make\_confined\_concrete()}).
    All functions are in the \texttt{fall\_n::} namespace.  \\
\addlinespace
\texttt{src/materials/RCSectionBuilder.hh}
  & Two designated-initialiser--friendly specification structs
    (\texttt{RCColumnSpec}, \texttt{RCBeamSpec}) and corresponding
    builder functions:
    \texttt{make\_rc\_column\_section(spec)} and
    \texttt{make\_rc\_beam\_section(spec)}, each returning a
    \texttt{Material<TimoshenkoBeam3D>}. \\
\addlinespace
\texttt{src/model/LoadUtilities.hh}
  & \texttt{apply\_uniform\_shell\_surface\_load()} template---integrates
    a uniform surface traction over shell element quadrature points.
    Extracted for reuse; not currently invoked by the building example
    (slab gravity uses a manual tributary approach). \\
\bottomrule
\end{tabular}
\end{center}

All four headers are registered in the umbrella include
\texttt{header\_files.hh} and follow the project's relative-path
include convention.

% ------------------------------------------------------------------
\subsection{Designated-Initialiser API}
\label{sec:phase2-designated-init}

The old factory functions captured constants from the enclosing
anonymous namespace, making them non-reusable.  The new API uses
C++20 designated initialisers to make every parameter explicit:

\begin{lstlisting}[language=C++,basicstyle=\small\ttfamily]
// Old: zero-argument, 60 lines, non-reusable
static Material<TimoshenkoBeam3D> make_column_material() { /* ... */ }

// New: all parameters in a self-documenting spec struct
auto col_mat = fall_n::make_rc_column_section({
    .b            = COLUMN_B,
    .h            = COLUMN_H,
    .cover        = COLUMN_COVER,
    .bar_diameter = COLUMN_BAR_D,
    .tie_spacing  = COLUMN_TIE_SPACING,
    .fpc          = COLUMN_FPC,
    .nu           = NU_RC,
    .steel_E      = STEEL_E,
    .steel_fy     = STEEL_FY,
    .steel_b      = STEEL_B,
    .tie_fy       = TIE_FY,
});
\end{lstlisting}

This pattern is applied to both the column and beam section builders.
The slab material (\texttt{MindlinShellMaterial}) remains a simple
three-parameter construction retained inline.

% ------------------------------------------------------------------
\subsection{Migration Statistics}
\label{sec:phase2-statistics}

\begin{itemize}
  \item \textbf{Lines removed from \texttt{main.cpp}}: $\sim$240
        (12 functions + associated includes).
  \item \textbf{Lines added to library}: $\sim$320 across four new
        headers (more thorough parameterisation and documentation).
  \item \textbf{Remaining local functions}: \texttt{node\_id()},
        \texttt{make\_slab\_material()},
        \texttt{add\_regular\_building\_nodes()},
        \texttt{add\_columns()}, \texttt{add\_beams()},
        \texttt{add\_slabs()}, \texttt{build\_structural\_elements()},
        and \texttt{run\_corotational\_dynamic\_rc\_building()}.
        These are model-construction and analysis-driver functions
        targeted for Phase~3 (\texttt{ModelBuilder}).
  \item \textbf{Unused constant removed}: \texttt{KAPPA\_RC}---the
        shear correction factor is now embedded in the
        \texttt{RCColumnSpec}/\texttt{RCBeamSpec} structs
        (default $\kappa_y = \kappa_z = 5/6$).
\end{itemize}

% ------------------------------------------------------------------
\subsection{Build and Test Verification}
\label{sec:phase2-verification}

After the extraction and refactoring:

\begin{center}
  \textbf{48 / 48 build targets compiled (zero warnings, \texttt{-Werror}).}\\[4pt]
  \textbf{45 / 45 CTests passed, 0 failures.}
\end{center}

No functional changes were introduced---the building example produces
identical results before and after Phase~2.


% ══════════════════════════════════════════════════════════════════════════
\section{Build System Optimisation (Phase~2B)}
\label{sec:build-optimisation}

As the test suite grew to 45 tests spanning from sub-second unit
checks to three-minute nonlinear dynamic analyses, the
\texttt{CMakeLists.txt} became a maintenance burden ($\sim$870 lines
of repetitive boilerplate) and full test execution time reached
$\sim$3.5~minutes.  This section documents the refactoring of the
build system for efficient development iteration.

% ------------------------------------------------------------------
\subsection{CMake Helper Function}
\label{sec:cmake-helper}

A single CMake function, \texttt{fall\_n\_add\_test()}, replaces all
per-test boilerplate.  It accepts five keyword arguments:

\begin{lstlisting}[language={},basicstyle=\small\ttfamily,
  morekeywords={fall_n_add_test,TARGET,NAME,SOURCE,LIBS,LABEL}]
fall_n_add_test(
    TARGET  fall_n_beam_test
    NAME    beam_element
    SOURCE  tests/test_beam_element.cpp
    LIBS    FULL          # FULL | CORE | EIGEN | VTK | NONE
    LABEL   unit          # unit | integration | slow
)
\end{lstlisting}

The \texttt{LIBS} tier controls which dependency bundles are linked:

\begin{center}
\small
\begin{tabular}{lp{8cm}}
\toprule
\textbf{Tier} & \textbf{Libraries} \\
\midrule
\texttt{FULL} & MPI + PETSc + Eigen + VTK (with \texttt{vtk\_module\_autoinit}) \\
\texttt{CORE} & MPI + PETSc + Eigen \\
\texttt{EIGEN}& Eigen only \\
\texttt{VTK}  & VTK only (with autoinit) \\
\texttt{NONE} & Header-only (include directory only) \\
\bottomrule
\end{tabular}
\end{center}

OpenMP is linked automatically on all tiers (except \texttt{NONE})
when \texttt{find\_package(OpenMP)} succeeds.

\paragraph{Result.}  \texttt{CMakeLists.txt} reduced from $\sim$870
lines of test declarations to $\sim$140 lines (one call per test,
three lines each).  Adding a new test requires exactly one
\texttt{fall\_n\_add\_test()} invocation.

% ------------------------------------------------------------------
\subsection{CTest Label Hierarchy}
\label{sec:ctest-labels}

Tests are classified into three labels based on measured wall-clock
time:

\begin{center}
\small
\begin{tabular}{lcrl}
\toprule
\textbf{Label} & \textbf{Count} & \textbf{Wall time} & \textbf{Use case} \\
\midrule
\texttt{unit}        & 36 & $\sim$1--6\,s   & Edit--compile--test inner loop \\
\texttt{integration} &  3 & $\sim$7--10\,s  & Pre-push gate \\
\texttt{slow}        &  6 & $\sim$2--4\,min & Nightly / CI full regression \\
\midrule
\textbf{Total}       & 45 & $\sim$3.5\,min  & Complete suite \\
\bottomrule
\end{tabular}
\end{center}

Slow tests also receive a 600\,s timeout guard via
\texttt{set\_tests\_properties(\ldots\ PROPERTIES TIMEOUT 600)}.

% ------------------------------------------------------------------
\subsection{CMake Presets}
\label{sec:cmake-presets}

A \texttt{CMakePresets.json} file provides three configure presets
and three test presets, enabling one-command workflows:

\begin{lstlisting}[language=bash,basicstyle=\small\ttfamily]
# Fast feedback (~6 s):
ctest --preset unit

# Pre-push gate (~10 s):
ctest --preset fast      # unit + integration

# Full regression (~3 min):
ctest --preset all
\end{lstlisting}

Configure presets include \texttt{debug} (default, current flags),
\texttt{release} (\texttt{-O2}), and \texttt{relwithdebinfo} (for
profiling).

% ------------------------------------------------------------------
\subsection{Compiler Cache}
\label{sec:ccache}

The build system now detects \texttt{ccache} via
\texttt{find\_program()} and sets
\texttt{CMAKE\_CXX\_COMPILER\_LAUNCHER} automatically.  On
incremental rebuilds after header-only changes, \texttt{ccache} can
reduce recompilation time by 50--80\% by serving cached object files.
Install with \texttt{sudo apt install ccache} to activate.

% ------------------------------------------------------------------
\subsection{Development Workflow}
\label{sec:dev-workflow}

The recommended workflow during active development is:

\begin{enumerate}
  \item Edit source files.
  \item \texttt{ninja} (incremental build, typically $<$5\,s for
        single-file changes).
  \item \texttt{ctest --preset unit} (36 unit tests, $\sim$6\,s).
  \item When satisfied: \texttt{ctest --preset fast} (39 tests,
        $\sim$10\,s).
  \item Before committing: \texttt{ctest --preset all} (full 45-test
        regression, $\sim$3\,min).
\end{enumerate}

This tiered approach reduces the inner-loop iteration time from
$\sim$3.5~minutes (running all 45 tests every time) to under
10~seconds, while still guaranteeing full coverage before each commit.


% ══════════════════════════════════════════════════════════════════════════
\section{Summary}
\label{sec:obs-summary}

The proposed architecture introduces two orthogonal subsystems:
\begin{itemize}
  \item An \textbf{Observer/Recorder} pipeline that replaces the
        monolithic monitor lambda with composable, single-responsibility
        observers---each independently testable and reconfigurable at
        runtime.  The architecture deliberately departs from OpenSees'
        pull/file/scripting model in favour of a push/in-memory/C++
        model (see \S\ref{sec:opensees-comparison}).
        \emph{Phase~1 is complete}: seven concrete observers,
        a dual static/dynamic composition model, and full integration
        into \texttt{DynamicAnalysis} with backward compatibility.
  \item \textbf{Helper extraction} (Phase~2) that promotes $\sim$240
        lines of anonymous-namespace utilities and fiber-section
        factories from \texttt{main.cpp} into four reusable library
        headers under the \texttt{fall\_n::} namespace, using C++20
        designated-initialiser specification structs for clean,
        self-documenting APIs.  \emph{Phase~2 is complete}: all 45
        CTests pass; the build compiles with zero warnings under
        \texttt{-Werror}.
  \item \textbf{Build system optimisation} (Phase~2B) that reduces
        \texttt{CMakeLists.txt} from $\sim$870 lines to $\sim$340 via
        a \texttt{fall\_n\_add\_test()} helper function, introduces
        CTest labels (\texttt{unit}/\texttt{integration}/\texttt{slow})
        for tiered test execution, and provides CMake presets for
        one-command workflows (\texttt{ctest --preset unit} finishes
        in $\sim$6\,s).
  \item A \textbf{ModelBuilder} orchestrator that extracts
        the domain/model-construction functions into spec-driven
        builders with fluent APIs, keeping the door open for
        future scripting-language bindings
        (Phase~3, see \S\ref{sec:phase3-implementation}).
\end{itemize}

Remaining phases: completing the
\texttt{std::print}/\texttt{format} migration across the wider
library (Phase~4).


% ══════════════════════════════════════════════════════════════════════════
\section{Phase~3 Implementation Notes: ModelBuilder}
\label{sec:phase3-implementation}

Phase~3 implements the ModelBuilder architecture proposed in
\S\ref{sec:model-builder}, extracting the \emph{remaining}
anonymous-namespace functions from \texttt{main.cpp}---grid node
generation, element creation, physical-group--to--material mapping,
self-weight computation, and lateral load application---into two
new library headers under the \texttt{fall\_n::} namespace.

% ------------------------------------------------------------------
\subsection{Design Decisions}
\label{sec:phase3-decisions}

\paragraph{Spec-driven, not OOP-driven.}
The ch54 design proposal (\S\ref{sec:model-builder}) offered a
\texttt{ModelBuilder<dim>} class with a fluent API.  The implementation
departs from this by choosing \emph{free functions} and \emph{plain
data structs} over a stateful class hierarchy.  The reasons:

\begin{enumerate}
  \item \textbf{Scripting-friendliness.}  A function that takes a
        POD-like struct and returns a result is the simplest possible
        FFI boundary.  Julia's \texttt{CxxWrap.jl} and Python's
        \texttt{pybind11} can expose \texttt{make\_building\_domain(spec)}
        with minimal glue code---no need to wrap a fluent builder class
        with chaining semantics.

  \item \textbf{Composability.}  Free functions compose naturally:
        \texttt{make\_building\_domain} $\to$
        \texttt{StructuralModelBuilder::build\_elements} $\to$
        Model constructor.  Each step is independently testable and
        replaceable (e.g.\ swap \texttt{make\_building\_domain} for a
        Gmsh import).

  \item \textbf{Minimum state.}  The only stateful class is
        \texttt{StructuralModelBuilder}, which accumulates
        group$\to$material mappings via a fluent API and then
        produces elements in a single \texttt{build\_elements()} call.
        This is a \emph{configuration object}, not a lifecycle manager.
\end{enumerate}

\paragraph{Non-template StructuralModelBuilder public API.}
While the builder is technically a class template (defaulting to
\texttt{BeamElement<TimoshenkoBeam3D, 3, beam::Corotational>} and
\texttt{ShellElement<MindlinReissnerShell3D>}), the default template
arguments mean it can be instantiated as simply
\texttt{fall\_n::StructuralModelBuilder\{\}}.  Users who need a
different kinematic policy (e.g.\ \texttt{beam::SmallRotation}) can
supply explicit template arguments, but the common case requires zero
template syntax---a deliberate choice for scripting-language
compatibility.

\paragraph{BuildingGrid: topology metadata as a return value.}
The \texttt{BuildingGrid} struct returned alongside the
\texttt{Domain<3>} captures the grid dimensions and provides O(1)
\texttt{node\_id(ix, iy, level)} lookups.  This avoids recalculating
grid indices when applying loads or querying results, and serves as
the bridge between the declarative spec and subsequent analysis steps.
The grid also exposes topology counters
(\texttt{num\_columns()}, \texttt{num\_beams()}, \texttt{num\_slabs()})
used in the final report.

\paragraph{Designated-initialiser API (continued from Phase~2).}
Both \texttt{FramedBuildingSpec} and \texttt{SelfWeightSpec} use C++20
aggregate initialisation with designated initialisers, consistent with
the \texttt{RCColumnSpec}/\texttt{RCBeamSpec} pattern established in
Phase~2.  This makes each parameter self-documenting at the call site.

% ------------------------------------------------------------------
\subsection{New File Organisation}
\label{sec:phase3-files}

\begin{center}
\small
\begin{tabular}{lp{8cm}}
\toprule
\textbf{Header} & \textbf{Contents} \\
\midrule
\texttt{src/model/BuildingDomainBuilder.hh}
  & \texttt{FramedBuildingSpec} (user-facing building specification),
    \texttt{BuildingGrid} (immutable topology metadata with
    \texttt{node\_id()} and topology counters),
    \texttt{make\_building\_domain(spec)} (free function returning
    \texttt{pair<Domain<3>, BuildingGrid>}),
    \texttt{SelfWeightSpec} + \texttt{apply\_building\_self\_weight()}
    (tributary gravity loads),
    \texttt{apply\_triangular\_lateral\_load()} (height-proportional
    lateral forces). \\
\addlinespace
\texttt{src/model/StructuralModelBuilder.hh}
  & \texttt{StructuralModelBuilder<FrameElemT, ShellElemT>}
    class template with fluent \texttt{set\_frame\_material()} /
    \texttt{set\_shell\_material()} and a single
    \texttt{build\_elements(domain)} that produces the type-erased
    \texttt{StructuralElement} container. \\
\bottomrule
\end{tabular}
\end{center}

Both headers are registered in the umbrella include
\texttt{header\_files.hh}.

% ------------------------------------------------------------------
\subsection{API Examples}
\label{sec:phase3-api}

\paragraph{Domain construction.}
\begin{lstlisting}[language=C++,basicstyle=\small\ttfamily]
auto [domain, grid] = fall_n::make_building_domain({
    .x_axes       = {0.0, 6.0, 12.0, 18.0},
    .y_axes       = {0.0, 5.0, 10.0},
    .num_stories  = 5,
    .story_height = 3.20,
});
\end{lstlisting}

\paragraph{Physical-group--to--material mapping.}
\begin{lstlisting}[language=C++,basicstyle=\small\ttfamily]
auto elements = fall_n::StructuralModelBuilder{}
    .set_frame_material("Columns", col_mat)
    .set_frame_material("Beams",   beam_mat)
    .set_shell_material("Slabs",   slab_mat)
    .build_elements(domain);

StructuralModel model{domain, std::move(elements)};
\end{lstlisting}

\paragraph{Self-weight and lateral loads.}
\begin{lstlisting}[language=C++,basicstyle=\small\ttfamily]
fall_n::apply_building_self_weight(model, grid, {
    .density        = RC_DENSITY,
    .gravity        = GRAVITY_ACCEL,
    .column_area    = COLUMN_AREA,
    .beam_area      = BEAM_AREA,
    .slab_thickness = SLAB_T,
});

fall_n::apply_triangular_lateral_load(model, grid,
    0 /* X direction */, LATERAL_AMPLITUDE);
\end{lstlisting}

% ------------------------------------------------------------------
\subsection{Refactoring Impact on \texttt{main.cpp}}
\label{sec:phase3-refactoring}

\begin{itemize}
  \item \textbf{Lines removed}: $\sim$240 (9~functions:
        \texttt{node\_id()}, 3~material factories, 4~geometry
        generators, \texttt{build\_structural\_elements()}, and the
        \texttt{Z\_LEVELS} compile-time lambda).
  \item \textbf{Lines added to library}: $\sim$290 across two new
        headers (more thorough parameterisation and reusability).
  \item \textbf{\texttt{main.cpp} size}: 628~$\to$~389 lines
        ($-$38\%).
  \item \textbf{Remaining local code}: constants (problem-specific
        data), type aliases, the analysis driver function
        (\texttt{run\_corotational\_dynamic\_rc\_building}) which now
        reads as a declarative sequence of builder calls, and the
        \texttt{main()} entry point.
\end{itemize}

\paragraph{Cumulative reduction.}
Across Phases~1--3, \texttt{main.cpp} has been reduced from its
pre-Phase~1 size of $\sim$870 lines to $\sim$389 lines ($-$55\%),
with all extracted logic promoted to reusable, namespace-qualified
library headers.

% ------------------------------------------------------------------
\subsection{Scripting-Language Openness}
\label{sec:phase3-scripting}

The \texttt{FramedBuildingSpec} and \texttt{SelfWeightSpec} structs
contain only \texttt{double}, \texttt{int}, and
\texttt{std::vector<double>} members---types that map directly to
native arrays and scalars in Python (\texttt{pybind11::list}) and
Julia (\texttt{Vector\{Float64\}}).

The free-function entry points (\texttt{make\_building\_domain},
\texttt{apply\_building\_self\_weight},
\texttt{apply\_triangular\_lateral\_load}) have simple signatures:
\texttt{(const Spec\&) -> Result}.  No template arguments appear at
call sites when using the default element types.

A future binding layer would expose these as:
\begin{lstlisting}[language=Python,basicstyle=\small\ttfamily]
# Python (pybind11)
domain, grid = fall_n.make_building_domain(
    x_axes=[0, 6, 12, 18], y_axes=[0, 5, 10],
    num_stories=5, story_height=3.2)
\end{lstlisting}

\begin{lstlisting}[language=Julia,basicstyle=\small\ttfamily]
# Julia (CxxWrap.jl)
domain, grid = FallN.make_building_domain(
    x_axes=[0, 6, 12, 18], y_axes=[0, 5, 10],
    num_stories=5, story_height=3.2)
\end{lstlisting}

The key design principle: \emph{no template syntax is required at the
scripting boundary}.  All templates are resolved at library
compilation time; the scripting layer only sees concrete types and
free functions.

% ------------------------------------------------------------------
\subsection{Build and Test Verification}
\label{sec:phase3-verification}

After the extraction and refactoring:

\begin{center}
  \textbf{Build: zero warnings (clang++-20, \texttt{-Werror}).} \\[4pt]
  \textbf{45 / 45 CTests passed, 0 failures.}
\end{center}

No functional changes---the building example produces identical
results before and after Phase~3.


% ══════════════════════════════════════════════════════════════════════════
\section{Scripting Language Evaluation: Julia vs.\ Python}
\label{sec:scripting-evaluation}

A design requirement stated from the outset of Phase~3 is that the
ModelBuilder API must remain open to a future scripting-language
binding.  This section presents a systematic comparison of the two
strongest candidates---\textbf{Julia} and \textbf{Python}---evaluated
against the specific characteristics and target audience of
\texttt{fall\_n}, and closes with a justified recommendation.

% ------------------------------------------------------------------
\subsection{Evaluation Criteria}
\label{sec:scripting-criteria}

The comparison is organised along seven axes relevant to a
computational structural mechanics library:

\begin{enumerate}
  \item \textbf{FFI maturity and C++ interoperability.}
        How seamless is the Foreign Function Interface for wrapping
        a modern C++23 library that relies on PETSc, Eigen, and VTK?

  \item \textbf{Numerical ecosystem.}
        Availability of performant linear algebra, sparse solvers,
        optimisation, and post-processing tools native to the
        language.

  \item \textbf{Structural engineering community penetration.}
        Existing FEM/structural analysis software with established
        scripting APIs in the language.

  \item \textbf{Interactive workflows.}
        Notebook support (Jupyter / Pluto), REPL-driven
        exploration, and integration with plotting libraries.

  \item \textbf{Performance profile.}
        Overhead of the FFI boundary, ability to write
        performance-critical callbacks (e.g.\ time-dependent force
        functions) without falling back to C++.

  \item \textbf{Distribution and deployment.}
        Ease of distributing a compiled C++ library with its
        scripting bindings to end users (package managers,
        virtual environments, reproducibility).

  \item \textbf{User base accessibility.}
        How many structural engineers, graduate students, and
        researchers already know or are likely to learn the language?
\end{enumerate}

% ------------------------------------------------------------------
\subsection{Python: Strengths and Limitations}
\label{sec:python-analysis}

\paragraph{FFI via pybind11 / nanobind.}
\texttt{pybind11}~\cite{pybind11} is the de-facto standard for
exposing C++ to Python.  It handles aggregate structs, STL containers,
\texttt{std::string}, and opaque pointer wrappers with minimal glue
code.  The newer \texttt{nanobind} offers improved compilation
speed.  Both produce standard CPython extensions distributable
via \texttt{pip}.

A key advantage: PETSc's own Python layer (\texttt{petsc4py})
and VTK's Python bindings are production-grade and installable
via \texttt{pip}.  This means a Python user can mix
\texttt{fall\_n} calls with direct PETSc vector/matrix queries
and VTK rendering---all in the same process.

\paragraph{Ecosystem dominance for engineering.}
\begin{itemize}
  \item \textbf{OpenSeesPy}~\cite{OpenSees}: The most widely used
        nonlinear structural FEM framework exposes its entire API
        through Python, making structural engineers already literate
        in Python-based FEM workflows.
  \item \textbf{SAP2000 / ETABS OAPI}: CSi products expose COM/Python
        scripting interfaces.
  \item \textbf{Abaqus scripting}: Python is the native scripting
        language for Abaqus CAE.
  \item Post-processing: NumPy, SciPy, Matplotlib, pandas, and
        Jupyter are ubiquitous in the engineering research workflow.
\end{itemize}

\paragraph{Limitations.}
\begin{itemize}
  \item \textbf{Performance ceiling.}  Pure-Python callbacks
        (e.g.\ a time-dependent force function called at every time
        step) incur interpreter overhead.  For the 5\,000-step
        dynamic analysis in the building example, this translates
        to $\sim$5\,000 Python$\to$C++ round trips---manageable
        but not zero-cost.
  \item \textbf{GIL constraints.}  The Global Interpreter Lock
        prevents true multi-threaded callbacks.  Not a concern
        for \texttt{fall\_n} (parallelism is handled by MPI/PETSc
        below the binding interface), but limits future
        observer-level parallelism in Python.
  \item \textbf{Type system.}  Python's dynamic typing offers
        no compile-time safety.  Typing errors surface at runtime
        rather than at import time.
\end{itemize}

% ------------------------------------------------------------------
\subsection{Julia: Strengths and Limitations}
\label{sec:julia-analysis}

\paragraph{FFI via CxxWrap.jl.}
\texttt{CxxWrap.jl}~\cite{CxxWrap} generates Julia-callable
wrappers from C++ code annotated with a registration macro.
It handles parametric types, inheritance, and value/reference
semantics.  Unlike pybind11, the C++ wrapper is compiled into a
shared library, and Julia loads it via \texttt{ccall}.

Julia also has a zero-overhead \texttt{ccall} mechanism for C
functions, but C++23 templates and class hierarchies require
\texttt{CxxWrap.jl} for ergonomic wrapping.

\paragraph{Performance.}
Julia's JIT-compiled functions achieve C-like performance.  A
time-dependent force callback written in Julia and passed to the
solver would be compiled to native code \emph{before} the time-stepping
loop begins---eliminating the per-call interpreter overhead that
Python would incur.  For tight inner loops (e.g.\ user-defined
material laws, custom load functions), this is a significant
practical advantage.

\paragraph{Numerical native design.}
Julia was designed \emph{ab initio} for numerical computing:
\begin{itemize}
  \item Multi-dimensional arrays are first-class, column-major
        (matching Fortran/Eigen conventions), with BLAS/LAPACK
        natively integrated.
  \item Multiple dispatch enables type-stable generic code---a
        user can define a new material as a Julia struct and
        dispatch on it seamlessly.
  \item The \texttt{DifferentialEquations.jl} ecosystem provides
        production-quality time-stepping algorithms; future
        integration with \texttt{fall\_n}'s PETSc solver could
        provide Julia-side adaptive control.
\end{itemize}

\paragraph{Emerging FEM ecosystem.}
\begin{itemize}
  \item \texttt{Ferrite.jl}: A Julia-native FEM assembly library
        with modern API design.
  \item \texttt{Gridap.jl}: High-level FE discretisation framework.
  \item \texttt{PETSc.jl}: Native Julia bindings to PETSc.
  \item The community is smaller but rapidly growing in
        computational mechanics.
\end{itemize}

\paragraph{Limitations.}
\begin{itemize}
  \item \textbf{Community size in structural engineering.}
        No major structural analysis package (OpenSees, SAP2000,
        ETABS, Abaqus) exposes a Julia API.  Structural engineers
        are overwhelmingly Python-trained.
  \item \textbf{CxxWrap.jl maturity.}
        While functional, \texttt{CxxWrap.jl} has fewer users and
        less documentation than pybind11.  Complex wrapping scenarios
        (e.g.\ type-erased \texttt{StructuralElement}) require more
        manual glue.
  \item \textbf{First-call latency (TTFX).}
        Julia's JIT compilation introduces a ``time to first
        execution'' penalty---annoying for interactive exploration,
        although mitigated by \texttt{PackageCompiler.jl} and the
        runtime improvements in Julia~1.9+.
  \item \textbf{Distribution complexity.}
        Distributing a Julia package with a compiled C++ backend
        requires the \texttt{BinaryBuilder.jl} / \texttt{JLL}
        infrastructure, which is more complex than Python's
        \texttt{pip install}.
\end{itemize}

% ------------------------------------------------------------------
\subsection{Head-to-Head Comparison}
\label{sec:scripting-comparison}

\begin{center}
\small
\begin{tabular}{p{3.5cm}p{5.5cm}p{5.5cm}}
\toprule
\textbf{Criterion} & \textbf{Python (pybind11)} & \textbf{Julia (CxxWrap.jl)} \\
\midrule
FFI maturity
  & Production-grade; thousands of C++ libraries wrapped
  & Functional; fewer complex-wrapping examples \\
\addlinespace
PETSc integration
  & \texttt{petsc4py} (official, pip-installable)
  & \texttt{PETSc.jl} (community, functional) \\
\addlinespace
VTK integration
  & Official Python bindings (pip)
  & \texttt{WriteVTK.jl} (pure Julia writer) \\
\addlinespace
Structural eng.\ community
  & Dominant (OpenSeesPy, SAP2000, Abaqus)
  & Minimal (no major FEM tool uses Julia) \\
\addlinespace
Callback performance
  & Interpreter overhead ($\sim\mu$s/call)
  & JIT-compiled ($\sim$ns/call, C-like) \\
\addlinespace
Numerical expressiveness
  & Good (NumPy), but duck-typed
  & Excellent (native arrays, multiple dispatch, type-stable) \\
\addlinespace
Interactive notebooks
  & Jupyter (mature)
  & Jupyter + Pluto.jl (mature) \\
\addlinespace
Distribution
  & \texttt{pip wheel} (simple, well-understood)
  & \texttt{BinaryBuilder.jl} (more complex) \\
\addlinespace
User onboarding
  & Most engineers know Python
  & Learning curve; smaller community \\
\bottomrule
\end{tabular}
\end{center}

% ------------------------------------------------------------------
\subsection{Recommendation: Julia as Primary, Python as Secondary}
\label{sec:scripting-recommendation}

Despite Python's larger user base and more mature FFI tooling,
this document recommends \textbf{Julia as the primary scripting
layer} for \texttt{fall\_n}.  The reasoning is as follows.

\paragraph{1.\ Performance alignment.}
\texttt{fall\_n} is a \emph{research-oriented computational mechanics
library}, not a commercial pre/post-processor.  Its target users write
custom material models, time-dependent force functions, and
analysis control loops.  In Julia these are all JIT-compiled to
native code; in Python they remain interpreted.  The performance
gap matters most precisely where \texttt{fall\_n} differentiates
itself: the inner loop of nonlinear time-stepping with user-defined
callbacks.

\paragraph{2.\ Type-system congruence.}
Julia's parametric type system and multiple dispatch mirror the
C++ template-based design of \texttt{fall\_n}.  A Julia
\texttt{struct RCColumnSpec} maps one-to-one to the C++ aggregate;
dispatch on material types in Julia can parallel the C++
\texttt{Material<Policy>} hierarchy without losing type information.
Python's duck typing degrades this structure to runtime checks.

\paragraph{3.\ Composability with the numerical ecosystem.}
Julia's \texttt{DifferentialEquations.jl} provides adaptive
time-stepping, sensitivity analysis, and uncertainty quantification
that could augment PETSc's TS solver.
\texttt{ForwardDiff.jl} enables automatic differentiation of
user-defined constitutive laws---a critical capability for
research into new material models.

\paragraph{4.\ The ``two-language problem'' avoidance.}
A Python binding inevitably creates a two-language workflow:
Python for scripting, C++ for performance-critical code.
Engineers who hit the Python performance ceiling must either
(a) learn C++ to add features, or (b) write Cython/numba hacks.
With Julia, the scripting layer \emph{is} the performance layer;
a prototype material model written in Julia's REPL can be
used directly in production analyses without rewriting.

\paragraph{5.\ Long-term trajectory.}
Julia's adoption in computational science is accelerating
(MIT, national labs, DOE-funded projects).  While Python is
dominant today, Julia's technical advantages for numerical
computing make it increasingly likely to become the standard
in research-oriented computational mechanics---the exact
audience of \texttt{fall\_n}.

\paragraph{Python as a secondary binding.}
Python support should \emph{not} be abandoned.  A thin
\texttt{pybind11} layer for the ModelBuilder API (the same
spec-driven free functions designed in Phase~3) provides:
\begin{itemize}
  \item Access for engineers who are reluctant to learn Julia.
  \item Integration with Jupyter notebooks and Matplotlib for
        publication-quality plots.
  \item A bridge to the OpenSeesPy community for comparative
        benchmarks.
\end{itemize}

The recommended implementation order is:

\begin{enumerate}
  \item \textbf{Phase~4a}: Julia binding via \texttt{CxxWrap.jl}
        for the core API (\texttt{make\_building\_domain},
        \texttt{StructuralModelBuilder}, analysis drivers).
        This establishes the primary interactive workflow.
  \item \textbf{Phase~4b}: Python binding via \texttt{pybind11}
        for the same API surface, enabling cross-language
        access and community adoption.
\end{enumerate}

Both bindings share the same C++ API surface designed in
Phase~3 --- the spec-driven free functions and simple structs
that require \emph{no template syntax at the scripting boundary}.

% ------------------------------------------------------------------
\subsection{Architectural Implications}
\label{sec:scripting-architecture}

The choice of Julia as primary binding influences the C++ API
design in the following ways:

\begin{enumerate}
  \item \textbf{Callback-friendly solver interface.}
        The \texttt{DynamicAnalysis::set\_force\_function()} must
        accept a type-erased callable
        (\texttt{std::function<void(double, Vec)>}) that can be
        bound to a Julia function via \texttt{CxxWrap.jl}'s
        \texttt{@cxxfunction} macro.  This is already the case.

  \item \textbf{Return types over output parameters.}
        Julia's FFI works best with functions that \emph{return}
        results rather than filling output pointers.
        \texttt{make\_building\_domain()} returns a
        \texttt{pair<Domain, BuildingGrid>}---natural for Julia's
        destructuring.

  \item \textbf{String-based group names.}
        Julia's \texttt{String} maps cleanly to C++
        \texttt{std::string} via \texttt{CxxWrap.jl}.  The
        physical-group system in \texttt{fall\_n} uses strings
        throughout, which is FFI-optimal.

  \item \textbf{Thin C bridge for PETSc interop.}
        PETSc's C API is directly callable from Julia via
        \texttt{ccall}.  Vectors (\texttt{Vec}) and matrices
        (\texttt{Mat}) are opaque pointers---trivially passable
        across the FFI boundary.  A Julia user can call
        \texttt{VecNorm} on a gravity vector obtained from the
        model without any wrapper overhead.
\end{enumerate}

% ======================================================================
%  MITC Shell Elements and Corotational Shell Formulation
% ======================================================================

\section{Higher-Order MITC Shell Elements}
\label{sec:mitc-shells}

This section documents the design and implementation of MITC (Mixed
Interpolation of Tensorial Components) shell elements in
\texttt{fall\_n}, covering the MITC4, MITC9, and MITC16 families, their
assumed-strain methodology, and the integration of a corotational
kinematic framework for geometrically nonlinear analysis.

\subsection{The Transverse Shear Locking Problem}
\label{sec:shear-locking}

Standard isoparametric Mindlin--Reissner shell elements suffer from
\emph{transverse shear locking} when the shell thickness $t$ becomes
small relative to its characteristic span~$L$.  The mechanism is
well understood: the standard displacement-based interpolation of the
transverse shear strains $\gamma_{\alpha 3}$ ($\alpha=1,2$) introduces
spurious constraints that stiffen the element response as
$t/L \to 0$.

Consider a Mindlin--Reissner shell with mid-surface parametrisation
$\xi, \eta \in [-1,1]^2$.  The \emph{covariant} transverse shear
strains are
\begin{equation}
  \label{eq:cov-shear}
  \gamma_{\alpha 3}
  = \frac{\partial w}{\partial \xi_\alpha}
    + \mathbf{e}_3 \cdot
      \frac{\partial \mathbf{x}}{\partial \xi_\alpha}
    - \theta_\beta \,
      \frac{\partial \mathbf{x}}{\partial \xi_\alpha}
      \cdot \mathbf{e}_\beta
  ,
  \qquad \alpha = 1,2
\end{equation}
where $w$ is the transverse displacement, $\theta_\alpha$ are the
section rotations, and $\mathbf{x}(\xi,\eta)$ the mid-surface
mapping.  In a standard bilinear element the interpolations of $w$
and $\theta_\alpha$ both live in $Q_1$, but
$\partial w/\partial \xi_\alpha$ and the rotation contribution
$\theta_\beta (\partial \mathbf{x}/\partial \xi_\alpha)$ produce
\emph{bilinear} terms that cannot cancel over the full element
domain, creating parasitic shear energy.

The MITC methodology eliminates locking by replacing the directly
interpolated shear strains with an \emph{assumed} strain field
$\tilde{\gamma}_{\alpha 3}$ constructed from values sampled at
carefully chosen \emph{tying points}, where the directly computed
shear strains are reliable~\cite{DvorkinBathe1984, BatheDvorkin1986}.

\subsection{MITC4: Bilinear Assumed-Strain Shell}
\label{sec:mitc4}

The MITC4 element uses four tying points for the transverse shear
interpolation, two for each covariant component:

\begin{itemize}
  \item $\gamma_{13}$ is sampled at $(\xi, \eta) \in \{(0, -1), (0, +1)\}$
        and interpolated linearly in~$\eta$.
  \item $\gamma_{23}$ is sampled at $(\xi, \eta) \in \{(-1, 0), (+1, 0)\}$
        and interpolated linearly in~$\xi$.
\end{itemize}

The assumed shear strain field is then
\begin{align}
  \tilde{\gamma}_{13}(\xi, \eta)
  &= \tfrac12(1-\eta)\,\gamma_{13}^{(0,-1)}
   + \tfrac12(1+\eta)\,\gamma_{13}^{(0,+1)}
  ,
  \label{eq:mitc4-g13}
  \\
  \tilde{\gamma}_{23}(\xi, \eta)
  &= \tfrac12(1-\xi)\,\gamma_{23}^{(-1,0)}
   + \tfrac12(1+\xi)\,\gamma_{23}^{(+1,0)}
  .
  \label{eq:mitc4-g23}
\end{align}

This replaces the two shear-strain rows of the standard~$\mathbf{B}$
matrix with the assumed-strain operator
$\tilde{\mathbf{B}}_{\text{shear}}$, while the 6~membrane and
bending rows remain unmodified.  The resulting MITC4 element passes
all standard patch tests and exhibits asymptotically correct thin-shell
behaviour.

In \texttt{fall\_n}, the MITC4 assumed-strain policy is implemented as
\texttt{mitc::MITC4}.  Its static method
\begin{lstlisting}[language=C++, basicstyle=\small\ttfamily]
template <int TotalDofs>
static void compute_assumed_shear(
    double xi, double eta,
    const ElementGeometry<3>& geometry,
    const Eigen::Matrix3d& R,
    int dofs_per_node,
    Eigen::Matrix<double, 2, TotalDofs>& B_sh);
\end{lstlisting}
builds a $2 \times N_\mathrm{dof}$ assumed-shear $\tilde{\mathbf{B}}_\mathrm{shear}$
matrix at an arbitrary evaluation point.  The direct shear
$\mathbf{B}$-matrix at each tying point is computed using the
projected in-plane Jacobian from the helper
\texttt{mitc::in\_plane\_jacobian()}.

\subsection{MITC9: Biquadratic Assumed-Strain Shell}
\label{sec:mitc9}

The 9-node biquadratic shell requires a richer tying-point scheme to
avoid locking.  Following Bucalem and Bathe~\cite{BucalemBathe1993},
the MITC9 element uses 12~tying points:

\begin{itemize}
  \item 6~points for~$\gamma_{13}$: at
        $(\xi,\eta) \in \{-1/\sqrt{3},\, +1/\sqrt{3}\}
         \times \{-1,\, 0,\, +1\}$.
  \item 6~points for~$\gamma_{23}$: at
        $(\xi,\eta) \in \{-1,\, 0,\, +1\}
         \times \{-1/\sqrt{3},\, +1/\sqrt{3}\}$.
\end{itemize}

The assumed shear strains are expressed via Lagrange interpolation
on the tying grids:
\begin{equation}
  \tilde{\gamma}_{13}(\xi,\eta)
  = \sum_{i=1}^{2}\sum_{j=1}^{3}
    L_i(\xi) \, M_j(\eta) \;
    \gamma_{13}^{(A_i, B_j)}
  ,
  \qquad
  \tilde{\gamma}_{23}(\xi,\eta)
  = \sum_{i=1}^{3}\sum_{j=1}^{2}
    M_i(\xi) \, L_j(\eta) \;
    \gamma_{23}^{(C_i, D_j)}
  ,
  \label{eq:mitc9-assumed}
\end{equation}
where $L_i$ are the two-point Lagrange basis functions on
$\{-1/\sqrt{3},\, +1/\sqrt{3}\}$ and $M_j$ the three-point basis
on $\{-1,\, 0,\, +1\}$.

In \texttt{fall\_n}, this is encoded in \texttt{mitc::MITC9}, which
precomputes the $6 + 6$ tying-point direct shear matrices and then
blends them via the Lagrange weights at each evaluation point.  The
resulting $2 \times 54$ assumed-shear operator replaces rows 7--8 of
the standard~$\mathbf{B}$ matrix.

The MITC9 element uses a
$\texttt{LagrangeElement<3, 3, 3>}$ geometry (9-node Lagrangian
quadrilateral) with a $3 \times 3$ Gauss--Legendre quadrature
(\texttt{GaussLegendreCellIntegrator<3, 3>}).  Its 54~DOFs
($6 \times 9$~nodes) enable accurate representation of bending-dominated
problems on coarser meshes than the MITC4.

\subsection{MITC16: Bicubic Assumed-Strain Shell}
\label{sec:mitc16}

The 16-node bicubic shell extends the MITC concept to even higher
polynomial order, providing $p$-refinement capability.  Following
~\cite{BucalemBathe1993}, the MITC16 element uses 24~tying points:

\begin{itemize}
  \item 12~points for~$\gamma_{13}$: at
        $(\xi,\eta) \in
         \bigl\{-\sqrt{3/5},\, 0,\, +\sqrt{3/5}\bigr\}
         \times \{-1,\, -1/3,\, +1/3,\, +1\}$.
  \item 12~points for~$\gamma_{23}$: at
        $(\xi,\eta) \in
         \{-1,\, -1/3,\, +1/3,\, +1\}
         \times \bigl\{-\sqrt{3/5},\, 0,\, +\sqrt{3/5}\bigr\}$.
\end{itemize}

The assumed strain interpolation mirrors \eqref{eq:mitc9-assumed}
with a $3 \times 4$ (and $4 \times 3$) tensor-product Lagrange basis.
The three-point basis uses the Gauss--Legendre nodes
$\{-\sqrt{3/5},\, 0,\, +\sqrt{3/5}\}$ while the four-point basis
uses the equidistant nodes $\{-1,\, -1/3,\, +1/3,\, +1\}$.

The MITC16 uses a $\texttt{LagrangeElement<3, 4, 4>}$ geometry
(16-node bi-cubic Lagrangian quadrilateral) with $4 \times 4$
Gauss--Legendre quadrature.  Its 96~DOFs ($6 \times 16$~nodes)
make it suitable for single-element benchmarks and high-accuracy
shell analysis.

\subsection{Unified \texttt{MITCShellElement} Template}
\label{sec:mitc-template}

All three MITC families are instantiated from a single class
template:
\begin{lstlisting}[language=C++, basicstyle=\small\ttfamily]
template <
    std::size_t NNodes,        // 4, 9, or 16
    typename ShellPolicy,      // MindlinReissnerShell3D
    typename MITCPolicy,       // mitc::MITC4 / MITC9 / MITC16
    typename KinematicPolicy   // shell::SmallRotation or shell::Corotational
        = shell::SmallRotation,
    typename AsmPolicy = assembly::DirectAssembly>
class MITCShellElement;
\end{lstlisting}

\noindent
The template parameters encode:
\begin{description}
  \item[\texttt{NNodes}]
    Number of element nodes.  The total DOFs are $6 \times N$.
    Supported: 4 (bilinear), 9 (biquadratic), 16~(bicubic).
  \item[\texttt{ShellPolicy}]
    Material policy defining the strain measure (8-component
    Mindlin--Reissner: 3~membrane, 3~bending, 2~shear) and the
    constitutive relation (block-diagonal $\mathbf{D}$ matrix).
  \item[\texttt{MITCPolicy}]
    Assumed-strain policy that provides
    \texttt{compute\_assumed\_shear()} for the chosen tying-point
    scheme.
  \item[\texttt{KinematicPolicy}]
    Small-rotation or corotational strategy (Sec.~\ref{sec:corotational}).
  \item[\texttt{AsmPolicy}]
    Assembly policy (direct or matrix-free).
\end{description}

Convenient type aliases simplify usage:
\begin{lstlisting}[language=C++, basicstyle=\small\ttfamily]
using MITC4Shell<>  = MITCShellElement<4,  ..., mitc::MITC4>;
using MITC9Shell<>  = MITCShellElement<9,  ..., mitc::MITC9>;
using MITC16Shell<> = MITCShellElement<16, ..., mitc::MITC16>;
using CorotationalMITC4Shell<> = MITCShellElement<4, ..., shell::Corotational>;
using CorotationalMITC9Shell<> = MITCShellElement<9, ..., shell::Corotational>;
using CorotationalMITC16Shell<>= MITCShellElement<16,..., shell::Corotational>;
\end{lstlisting}

All aliases satisfy the \texttt{FiniteElement} concept verified via
\texttt{static\_assert} at the end of the header.

\subsection{Stiffness Matrix Assembly}
\label{sec:mitc-stiffness}

The element stiffness matrix is assembled in local coordinates as
\begin{equation}
  \mathbf{K}_{\mathrm{loc}}
  = \int_{\Omega_{\Box}} \mathbf{B}^T \, \mathbf{D} \, \mathbf{B}
    \; \lvert J \rvert \, d\xi\, d\eta
  ,
  \label{eq:K-local}
\end{equation}
where the $8 \times N_\mathrm{dof}$ strain--displacement matrix is
partitioned as
\[
  \mathbf{B} =
  \begin{pmatrix}
    \mathbf{B}_{\mathrm{membrane}} \\
    \mathbf{B}_{\mathrm{bending}}  \\
    \tilde{\mathbf{B}}_{\mathrm{shear}}
  \end{pmatrix}
  .
\]
The first six rows ($\mathbf{B}_\mathrm{membrane}$ and
$\mathbf{B}_\mathrm{bending}$) are computed from the standard
displacement interpolation; the last two rows
($\tilde{\mathbf{B}}_\mathrm{shear}$) come from the MITC
assumed-strain policy.  The material matrix $\mathbf{D}$ is
block-diagonal:
\[
  \mathbf{D} =
  \begin{pmatrix}
    \mathbf{D}_m  &                & \\
                  & \mathbf{D}_b   & \\
                  &                & \mathbf{D}_s
  \end{pmatrix}
  ,
\]
with membrane $\mathbf{D}_m = \frac{Et}{1-\nu^2}\mathbf{C}$,
bending $\mathbf{D}_b = \frac{Et^3}{12(1-\nu^2)}\mathbf{C}$, and
shear $\mathbf{D}_s = \kappa G t \, \mathbf{I}_2$ where
$\kappa = 5/6$ is the shear correction factor.

\subsection{Consistent Mass Matrix}
\label{sec:mitc-mass}

The consistent mass matrix is assembled as
\begin{equation}
  \mathbf{M} = \int_{\Omega_{\Box}}
    \rho \, t \;
    \mathbf{N}^T \mathbf{N}
    \; \lvert J \rvert \, d\xi\, d\eta
  ,
  \label{eq:M-consistent}
\end{equation}
where $\mathbf{N}$ is the $3 \times N_\mathrm{dof}$ shape function
matrix that picks out the translational DOFs.  This yields the
exact total mass $\rho \, t \, A$ when summed over all translational
DOF pairs:
$\sum_{a,b} M_{a\cdot6+k,\, b\cdot6+k} = \rho \, t \, A$
for each translational direction $k \in \{0,1,2\}$.
Rotational inertia terms are omitted (standard practice for
structural-dynamics codes), zeroing the $\theta$-DOF mass
subblocks.

\subsection{Internal Force Vector}
\label{sec:mitc-internal-force}

The element internal force vector in local coordinates is
\begin{equation}
  \mathbf{f}_{\mathrm{int,loc}}
  = \int_{\Omega_{\Box}} \mathbf{B}^T \, \boldsymbol{\sigma}
    \; \lvert J \rvert \, d\xi\, d\eta
  ,
  \label{eq:f-int}
\end{equation}
where $\boldsymbol{\sigma}$ is the 8-component generalised stress
vector (membrane forces, bending moments, transverse shears).  For
elastic material, $\boldsymbol{\sigma} = \mathbf{D}\,\mathbf{B}\,\mathbf{u}_\mathrm{loc}$.
For nonlinear material, $\boldsymbol{\sigma}$ is evaluated from the
section constitutive state at each Gauss point.

% ======================================================================
\section{Corotational Kinematic Policy for Shells}
\label{sec:corotational}

\subsection{Motivation: Geometric Compatibility in Mixed Models}
\label{sec:corotational-motivation}

The \texttt{fall\_n} dynamic RC building model combines corotational
Timoshenko beams (columns and beams) with shell slabs.  When the
beam elements use the \texttt{beam::Corotational} kinematic policy,
they filter out rigid-body rotations and work with deformational
(small-strain) local DOFs.  If the shells remain in a small-rotation
formulation, the kinematic mismatch at beam--shell junctions causes
\emph{parasitic restoring forces} at shared nodes---the shell ``sees''
the large nodal rotations as genuine deformation while the beam has
already subtracted the rigid motion.

This manifests as spurious oscillations and eventual instability in
dynamic analysis, particularly for slender slabs where bending
stiffness is small compared to the parasitic forces.  The solution
is to apply the \emph{same corotational treatment} to the shell
elements.

\subsection{Corotational Kinematics: Theory}
\label{sec:cr-theory}

The corotational approach decomposes the element motion into a
rigid-body part and a deformational part.  For an element with
nodes at current positions $\mathbf{x}_I$ and reference positions
$\mathbf{X}_I$:

\begin{enumerate}
  \item \textbf{Element frame.} An orthonormal frame
        $\mathbf{R} \in SO(3)$ is attached to the element and
        updated from the deformed configuration via the tangent
        vectors at the element centroid:
        \[
          \mathbf{g}_\alpha = \sum_I \frac{\partial N_I}{\partial \xi_\alpha}
          \bigg|_{\boldsymbol{\xi}=\mathbf{0}} \mathbf{x}_I
          ,\qquad \alpha = 1,2
          .
        \]
        A Gram--Schmidt orthonormalisation produces
        $\mathbf{e}_1, \mathbf{e}_2, \mathbf{e}_3 = \mathbf{e}_1 \times \mathbf{e}_2$.

  \item \textbf{Deformational displacements.} The local
        (deformational) translation at each node is obtained by
        subtracting the rigid-body displacement:
        \begin{equation}
          \mathbf{d}_I =
          \mathbf{R}_{\mathrm{curr}}
            \bigl(\mathbf{x}_I - \bar{\mathbf{x}}_{\mathrm{def}}\bigr)
          -
          \mathbf{R}_{\mathrm{ref}}
            \bigl(\mathbf{X}_I - \bar{\mathbf{X}}_{\mathrm{ref}}\bigr)
          ,
          \label{eq:cr-disp}
        \end{equation}
        where $\bar{\mathbf{x}}_{\mathrm{def}}$ and
        $\bar{\mathbf{X}}_{\mathrm{ref}}$ are the deformed and
        reference centroids, and $\mathbf{R}_{\mathrm{curr}}$,
        $\mathbf{R}_{\mathrm{ref}}$ are the current and reference
        element frames.

  \item \textbf{Deformational rotations.} The nodal rotation
        (drilling) DOF is reduced by the element rotation:
        \begin{equation}
          \tilde{\theta}_{z,I} = \theta_{z,I}
            - \Delta\alpha_{\mathrm{elem}}
          ,
          \label{eq:cr-drill}
        \end{equation}
        where $\Delta\alpha_{\mathrm{elem}}$ is the in-plane
        rotation angle of $\mathbf{R}_{\mathrm{curr}}$ relative
        to $\mathbf{R}_{\mathrm{ref}}$.  The transverse rotations
        ($\theta_x$, $\theta_y$) remain as-is (they describe
        section tilting, not rigid rotation).
\end{enumerate}

The stiffness assembly proceeds identically to the small-rotation
case, but using the deformational DOFs.  The resulting local
tangent stiffness $\mathbf{K}_\mathrm{loc}$ is then transformed
to global coordinates via the current frame's transformation
matrix $\mathbf{T}$:
\begin{equation}
  \mathbf{K}_{\mathrm{glob}} =
  \mathbf{T}^T \bigl(\mathbf{K}_\mathrm{loc}
  + \mathbf{K}_\sigma\bigr) \mathbf{T}
  ,
  \label{eq:K-glob}
\end{equation}
where $\mathbf{K}_\sigma$ is the geometric stiffness contribution
from the current stress state~\cite{FelippaHaugen2005}.

\subsection{Geometric Stiffness for Shells}
\label{sec:geom-stiffness}

The corotational formulation requires a geometric stiffness matrix
to account for the change of the element frame under load.  For
the shell, this takes the form
\begin{equation}
  (K_\sigma)_{IJ} = \sigma_{33}
  \int_{\Omega_{\Box}}
    \nabla N_I \cdot \nabla N_J
    \; \lvert J \rvert \, d\xi\, d\eta
  ,
  \label{eq:K-sigma}
\end{equation}
acting on the transverse displacement DOFs only.  Here $\sigma_{33}$
is the transverse normal stress resultant (approximately the
membrane force component normal to the shell surface).  In
\texttt{fall\_n}, the \texttt{shell::Corotational::geometric\_stiffness()}
method assembles this contribution at each Gauss point and adds it
to the local stiffness before the coordinate transformation.

\subsection{Implementation in \texttt{fall\_n}}
\label{sec:cr-implementation}

The corotational shell kinematic policy is encoded in
\texttt{shell::Corotational}, defined in
\texttt{ShellKinematicPolicy.hh}.  Key methods:

\begin{description}
  \item[\texttt{update\_frame()}]
    Recomputes $\mathbf{R}_{\mathrm{curr}}$ from deformed nodal
    positions using tangent vectors at the centroid and
    Gram--Schmidt orthonormalisation.

  \item[\texttt{extract\_local\_dofs()}]
    Two overloads:
    \begin{enumerate}
      \item \textbf{5-argument (geometry-aware):} Implements
            Eq.~\eqref{eq:cr-disp}--\eqref{eq:cr-drill} using
            centroid-relative deformational displacements and
            element-rotation subtraction from the drilling DOF.
            This overload requires access to the element geometry
            for nodal coordinates.
      \item \textbf{2-argument (fallback):} Simple
            $\mathbf{u}_\mathrm{loc} = \mathbf{T}\,\mathbf{u}_\mathrm{glob}$
            when geometry is not available.
    \end{enumerate}

  \item[\texttt{geometric\_stiffness()}]
    Assembles the $\nabla N_I \cdot \nabla N_J$ contribution
    from Eq.~\eqref{eq:K-sigma} on the $w$-DOFs using the current
    stress state.
\end{description}

In \texttt{MITCShellElement}, a private helper method
\texttt{extract\_local()} dispatches to the appropriate overload
via \texttt{if constexpr}:
\begin{lstlisting}[language=C++, basicstyle=\small\ttfamily]
auto extract_local(const FVectorT& u_glob,
                   const TMatrixT& T) const {
    if constexpr (std::same_as<KinematicPolicy,
                               shell::Corotational>) {
        return KinematicPolicy::extract_local_dofs(
            u_glob, T, R_, R_ref_, geometry_);
    } else {
        return KinematicPolicy::extract_local_dofs(u_glob, T);
    }
}
\end{lstlisting}

\subsection{Verification: Rigid-Body Invariance}
\label{sec:cr-verification}

A fundamental requirement of any corotational formulation is that
a pure rigid-body motion produces zero internal forces.  The test
\texttt{MITC4\_corotational\_rigid\_rotation\_zero\_force} applies
a $30^\circ$ in-plane rotation to all four nodes of a unit-square
MITC4 shell and verifies:
\[
  \lVert \mathbf{f}_\mathrm{int} \rVert_2 < 10^{-10}
  \quad\text{(achieved: } 4.49 \times 10^{-12}\text{)}
  .
\]
This confirms that the centroid-relative deformational displacement
extraction (Eq.~\ref{eq:cr-disp}) correctly removes the rigid
component to machine precision.

Additional verification tests include:
\begin{itemize}
  \item Corotational MITC4 at zero displacement gives identical
        stiffness to the small-rotation MITC4.
  \item Corotational MITC9 at zero displacement matches the
        small-rotation MITC9 (verified for 9-node elements).
  \item Stiffness symmetry and positive semi-definiteness
        (exactly 6 zero eigenvalues for an unconstrained shell
        patch) for all MITC families.
  \item $\mathbf{K}\,\mathbf{u} = \mathbf{f}$ consistency:
        applying the stiffness matrix to a generic displacement
        vector and comparing against the internal force vector,
        achieving residuals of $O(10^{-16})$.
\end{itemize}

\subsection{VTK Visualisation of Higher-Order Shells}
\label{sec:vtk-shells}

The VTK exporter has been generalised to handle $Q_4$, $Q_9$,
and $Q_{16}$ shell elements.  A helper structure
\texttt{detail::ShellGrid} generates the parametric sampling
grid and sub-quad connectivity for each element order:

\begin{center}
\begin{tabular}{l c c c}
  \toprule
  Element  & Nodes & Parametric grid & Sub-quads \\
  \midrule
  MITC4    & 4     & $2 \times 2$    & 1         \\
  MITC9    & 9     & $3 \times 3$    & 4         \\
  MITC16   & 16    & $4 \times 4$    & 9         \\
  \bottomrule
\end{tabular}
\end{center}

The \texttt{make\_shell\_grid(n\_nodes)} factory produces the
grid structure containing:
\begin{itemize}
  \item \texttt{points}: parametric coordinates for nodal sampling.
  \item \texttt{quads}: sub-quad connectivity for surface rendering.
  \item \texttt{edges}: boundary edges for axis (wireframe)
        rendering.
\end{itemize}

The \texttt{dispatch\_supported\_structural()} function has been
extended to dispatch all six MITC shell type variants
(MITC4/9/16 in both small-rotation and corotational modes).  The
type-erased \texttt{StructuralElement} wrapper's \texttt{as<T>()}
method is used for compile-time dispatch within the visitor pattern.

\subsection{Application: Corotational Dynamic RC Building}
\label{sec:cr-application}

The \texttt{main.cpp} dynamic 5-story RC building example has been
updated to use \texttt{CorotationalMITC4Shell<>} for the elastic
slab elements, specified via the \texttt{StructuralModelBuilder}
template parameter:

\begin{lstlisting}[language=C++, basicstyle=\small\ttfamily]
auto elements = fall_n::StructuralModelBuilder<
    BeamElement<TimoshenkoBeam3D, 3, beam::Corotational>,
    CorotationalMITC4Shell<>,
    TimoshenkoBeam3D,
    MindlinReissnerShell3D>{}
  .set_frame_material("Columns", column_material)
  .set_frame_material("Beams",   beam_material)
  .set_shell_material("Slabs",   slab_material)
  .build_elements(domain, &shell_geometries);
\end{lstlisting}

This ensures kinematic compatibility at beam--shell junctions:
both element types apply the corotational filter, preventing the
parasitic restoring forces that caused instability in the previous
formulation where shells used small-rotation kinematics alongside
corotational beams.

\subsection{Test Suite Summary}
\label{sec:mitc-tests}

The MITC shell element test suite (\texttt{test\_mitc\_shell\_elements.cpp})
contains 12~tests registered in CTest under the \texttt{unit} label:

\begin{center}
\small
\begin{tabular}{r l l}
  \toprule
  \# & Test name & Verification target \\
  \midrule
  1  & MITC4 matches original ShellElement     & Backward compatibility \\
  2  & MITC9 stiffness: symmetric, PSD, 6~zero & Stiffness properties \\
  3  & MITC16 stiffness: symmetric, PSD, 6~zero & Stiffness properties \\
  4  & CR-MITC4 zero at zero displacement       & Corotational consistency \\
  5  & CR-MITC9 matches SR at zero              & Corotational consistency \\
  6  & MITC9 vs MITC4 trace comparison          & $h$-convergence indicator \\
  7  & MITC9 $\mathbf{K}\mathbf{u}=\mathbf{f}$ & Force--stiffness consistency \\
  8  & MITC16 $\mathbf{K}\mathbf{u}=\mathbf{f}$ & Force--stiffness consistency \\
  9  & MITC9 GP strains finite                  & Assumed-strain correctness \\
  10 & MITC16 GP strains finite                 & Assumed-strain correctness \\
  11 & MITC9 consistent mass                    & Mass conservation \\
  12 & CR-MITC4 rigid rotation zero force       & Rigid-body invariance \\
  \bottomrule
\end{tabular}
\end{center}

All 12~tests pass, bringing the total test count to 46 (37~unit,
3~integration, 6~slow).
